\chapter{视频复原、增强与压缩}
\label{chap:vision-video-restoration}

一段夜间行车录像中，路牌在第$t$帧被运动模糊和噪声淹没，在第$t-3$帧和第$t+2$帧却各自清楚一些。我们希望得到一段能读出路牌的视频，可能还要把它放大四倍、补成慢动作，最后在有限带宽下传回监控中心。逐帧使用\ref{chap:vision-image-restoration}章的方法，会丢掉邻帧里现成的观测；直接把邻帧叠加到当前帧，又会把移动的车身纹理叠到路牌上，造成比单帧更难察觉的错误。

本章要回答的问题因此是：在规定的等待和成本之内，怎样从其他时刻取得当前画面真正需要的信息，怎样把它变成目标网格上的画面，以及怎样把相邻帧的共同内容变成解码端能够重现的少量比特。三项任务共用同一个前提，即先找到不同时刻画面之间的对应。\ref{chap:vision-video-foundations}章讲过块匹配和光流这两类经典运动估计，\ref{chap:vision-video-recognition}章又用RAFT估计像素位移；本章把这种对应用于读取像素和特征。复原与超分辨率利用对应“取回”邻帧内容，插帧利用对应“推算”未拍到的时刻，压缩则利用对应“少传”可预测的内容；它们对对应错误的承受方式不同，评价时要分别核验。

阅读本章需要\ref{chap:vision-image-restoration}章的退化模型、保真指标和图像压缩，以及\ref{chap:vision-video-foundations}章的运动估计与时间连续性度量。对应错误、遮挡、闪烁、等待和比特分配这几类困难并不是三项任务各自独有的，它们以不同强度同时出现，所以本章先把输入、输出和评价指标写成可检查的规定，再逐个说明这些挑战。后续各节的顺序跟着数据的处理过程走：先从邻帧取得并汇总可用信息，经典的运动补偿时域滤波和VBM4D与学习式的RVRT在这里对照；再把汇总结果落到原网格、更高分辨率网格或中间时刻，分别由RVRT、BasicVSR和FILM承担；编码要把同样的对应变成比特，先讲H.264/HEVC的混合编码框架，再看DVC把其中哪些模块换成网络并联合训练；最后回到评价，统一比较协议并同时核验内容、时间连续性、码率与延迟。图像压缩中的量化、概率模型和熵编码已在\ref{sec:vision-restoration-compression}节讲过，本章只说明视频新增的部分。

\section{视频复原、增强与压缩的问题与挑战}
\label{sec:vision-video-restoration-challenges}
\label{sec:vision-video-restoration-goals}

“让视频更清楚”可能指去噪、提亮、放大、补帧或减少压缩损伤。这些目标允许改变的量不同，所需的输入和参照也不同；同一个多帧网络在一种目标下合格，在另一种目标下可能越界。所以在调用任何模型之前，要先用可检查的方式写下三件事：输出在空间和时间上落在什么网格上，输出某一帧时允许读取哪些帧，以及用什么指标、在什么成本下判断结果合格。这三项规定确定之后，才能说清各种方法在什么地方困难，后面各节的模型也才能在同一口径下比较。

沿用\ref{chap:vision-image-restoration}章的记号，本章以$\bm{Y}_t$表示第$t$帧低质量观测，$\bm{X}_t$表示希望得到的目标帧，$\widehat{\bm{X}}_t$表示算法输出；$t=1,\dots,n$是帧序号，$n$为序列帧数，帧率记为$f$。每帧的高、宽和通道数分别为$H$、$W$和$c$。

\label{sec:vision-video-degradation}
规定输入与输出要回答两个问题：哪些跨帧变化属于应去除的退化，输出与输入在空间和时间网格上是什么关系。前者决定训练配对怎样构造，后者决定哪些输出位置有直接观测、哪些只能估计。

单帧中的噪声、模糊、低照度、分辨率和编码损失在视频中仍然存在，但可能随时间变化。自动曝光使同一路牌亮度改变，运动使车身边缘移动，遮挡使一块背景暂时不可见，切镜则让整个场景更换。只有任务约定要修正的变化才是退化；正常车辆运动不应被算法抹平，切镜也不应被“平滑”掉。

可以写$\bm{Y}_t=\mathcal{D}_t(\bm{X}_t)+\bm{\varepsilon}_t$概括第$t$帧的成像与处理，其中$\mathcal{D}_t$包括时变模糊、采样或编码等作用，$\bm{\varepsilon}_t$表示额外噪声。这只是任务层面的简化：长曝光会积分一段时间内的运动，不能总当作瞬时清晰图像上的固定卷积。多帧方法还要建立$\bm{X}_t$与其他时刻内容的对应，不能默认所有目标帧相同。

夜间录像中，平坦区域的随机颗粒提示噪声，移动路灯的拖尾可能来自曝光期间的运动，整帧明暗跳变可能来自曝光控制，方块边界则可能来自编码。不同成因需要不同的退化配对或数据标定。仅用合成高斯噪声训练的模型，不能据此宣称同时恢复了这些变化。

网格关系把本章的任务分成四类，表\ref{tab:vision-video-task-targets}列出它们的输入、输出与目标。复原和增强保持网格；超分辨率在每帧内增加像素，倍率记为$s$；插帧在两帧之间插入时刻$\tau\in(0,1)$；压缩的目标不是理想清晰帧，而是在比特预算内重现实际输入$\bm{Y}_t$。

\begin{table}[htbp]
  \centering
  \begin{tabularx}{\linewidth}{@{}lXXX@{}}
    \toprule
    任务 & 输入 & 输出网格 & 目标与直接观测\\
    \midrule
    复原与增强 & $\bm{Y}_{1},\dots,\bm{Y}_{n}$，每帧$H\times W$ & 同为$n$帧$H\times W$ & 清晰帧$\bm{X}_t$；每个输出位置都有退化观测\\
    超分辨率 & $n$帧$h\times w$ & $n$帧$sh\times sw$ & 高分辨率帧；新增位置只有邻帧的间接观测\\
    插帧 & 第$t$与$t+1$帧 & 时刻$t+\tau$的一帧 & 中间时刻画面；该时刻没有任何直接观测\\
    压缩 & $n$帧原视频 & 码流与重建帧$\widehat{\bm{Y}}_t$ & 重现$\bm{Y}_t$；解码端只见码流\\
    \bottomrule
  \end{tabularx}
  \caption{视频处理任务的网格关系。右列说明输出位置有无直接观测，它决定输出中哪些内容可以回查原始记录，哪些只是估计。}
  \label{tab:vision-video-task-targets}
\end{table}

\label{sec:vision-video-content-cost}
评价指标要与用途对应。观看可以容忍一定色调变化，车牌读取要求字符保持，速度测量还要求时间戳和几何关系可追溯。模型生成一个看起来清晰的字符，并不说明该字符存在于观测中；超分新增的高频细节、插帧得到的中间时刻，都需要与原始观测区分。保存原视频、处理结果及每帧所用的帧范围，才能事后检查一次改变来自输入证据还是模型先验。若用户要求把红车改成蓝车或改变行驶路线，问题已转为内容创作，见\ref{chap:vision-multimodal-generation}章。在这个前提下，本章使用四类可计算的指标：逐帧质量、时间一致性、码率和延迟。

逐帧质量沿用图像指标，只是要说明怎样在帧之间汇总。峰值信噪比（PSNR）、结构相似性（SSIM）和学习感知图像块相似度（LPIPS）已在\ref{sec:vision-restoration-challenges}节和\ref{sec:vision-restoration-psnr-ssim}节定义：PSNR由均方误差按式\eqref{eq:vision-psnr}换算，单位为分贝；SSIM在局部窗口中比较亮度、对比和结构；LPIPS比较预训练网络的多层特征。视频中，记第$t$帧输出与目标的均方误差为$\mathrm{MSE}_t$，相应峰值信噪比为$\mathrm{PSNR}_t=10\log_{10}(v_{\max}^2/\mathrm{MSE}_t)$，$v_{\max}$在8位图像中为255。视频论文通常先逐帧计算PSNR，再对帧取平均，DVC的补充材料就是这样说明的\scite{vision-dvc2019}。这与先把全部帧的均方误差平均、再换算成PSNR并不等价。教学数值：两帧的均方误差分别为1和100，逐帧PSNR为48.13和28.13 dB，平均38.13 dB；若先平均均方误差得到50.5，再换算则为31.10 dB。差帧的大均方误差在合并平均中占主导，逐帧平均对它的惩罚则轻得多，所以比较时要写明汇总方式。SSIM和LPIPS同样逐帧计算再平均。

逐帧指标看不到闪烁和字符随时间变化。时间一致性、闪烁和扭曲误差已在\ref{sec:vision-videofound-temporal-consistency}节定义（闪烁见式\eqref{eq:vision-videofound-flicker}，扭曲误差见式\eqref{eq:vision-videofound-warping-error}），本章使用其中的扭曲误差，并补充一种比较运动本身的指标。扭曲误差把相邻输出帧按运动对应搬到同一网格，在没有被遮挡的位置比较二者\scite{vision-lai2018blind}。按本章反向取样的约定，写成
\begin{equation}
  E_{\mathrm{warp}}=\frac{1}{n-1}\sum_{t=1}^{n-1}
  \frac{\sum_{\bm{p}}M_t(\bm{p})\,
    \bigl\lVert\widehat{\bm{X}}_t(\bm{p})
      -\mathcal{W}(\widehat{\bm{X}}_{t+1},\bm{u}_{t\to t+1})(\bm{p})\bigr\rVert_2^2}
  {\sum_{\bm{p}}M_t(\bm{p})}\text{，}
  \label{eq:vision-video-temporal-error}
\end{equation}
其中$\mathcal{W}(\cdot,\bm{u})$表示按位移场$\bm{u}$反向取样（\ref{sec:vision-video-warping}节给出计算），$\bm{u}_{t\to t+1}$是第$t$帧指向第$t+1$帧的光流，$M_t(\bm{p})\in\{0,1\}$是非遮挡掩码，取1表示位置$\bm{p}$在两帧中都可见。分母非零时式子才有定义。它就是式\eqref{eq:vision-videofound-warping-error}在整段上的平均，只是改用本章的输出记号与反向取样算子；若改用逐通道绝对误差，即得常见的L1变体，报告时应注明所用范数。这里有一个容易忽略的条件：对应和掩码必须来自与被测方法无关的估计，例如在原始输入上估计光流，否则方法自己的对齐错误会被同一份光流掩盖。教学数值：某可见位置当前输出为0.50，下一帧搬回来的值为0.46，平方差为0.0016；若下一帧在该处闪亮到0.60，平方差变为0.0196，扭曲误差随之增大。

扭曲误差有一个明显的漏洞：每帧都很模糊的输出，相邻帧之间也很接近。TecoGAN的作者指出，基于搬运后帧差的时间指标会被这类模糊结果欺骗，因而改为比较输出与真值“怎样运动”\scite{vision-tecogan2020}。他们定义
\begin{equation}
  \mathrm{tOF}_t=\bigl\lVert\mathrm{OF}(\widehat{\bm{X}}_{t-1},\widehat{\bm{X}}_t)
     -\mathrm{OF}(\bm{X}_{t-1},\bm{X}_t)\bigr\rVert_1\text{，}
  \label{eq:vision-video-tof}
\end{equation}
其中$\mathrm{OF}$是用固定的Farnebäck方法估计的光流，$\lVert\cdot\rVert_1$对全部位置和两个分量求绝对值之和；同文的$\mathrm{tLP}$把光流换成相邻帧之间的LPIPS距离。tOF需要真值视频：输出的运动与真值的运动越接近，值越小；模糊输出的光流往往与真值不同，就不能靠模糊得到低分。例如某位置真值相邻帧位移为$(1.5,0)$，输出估计为$(1.0,0)$，该位置贡献0.5。汇总时是求和还是按像素平均，应随结果写明。

码率衡量传输成本。对帧宽$W$、高$H$、共$n$帧、帧率$f$的视频，完整码流为$B$字节时，
\begin{equation}
  \text{平均比特/像素}=\frac{8B}{nWH},
  \qquad
  \text{平均比特率}=\frac{8Bf}{n}\ \mathrm{bit/s}\text{。}
  \label{eq:vision-video-bitrate}
\end{equation}
例如10秒、30帧每秒、$640\times360$的视频共300帧，完整码流恰为300000字节，则平均码率为240000 bit/s，平均约0.0347 bit/像素。这是构造的成本计算，不对应任何模型的实测表现。分母只计像素数，不乘颜色通道数。

一个编码器在不同参数下给出一条码率—质量曲线，单个工作点无法比较两条曲线。比较曲线常用BD码率（BD-rate），其定义见式\eqref{eq:vision-bd-rate}：对两个编码器各取若干工作点，用三次多项式把对数码率拟合成PSNR的函数，在两条曲线共同覆盖的质量区间上求对数码率的平均差，再换回比例\scite{vision-bjontegaard2001}。原提案取四个工作点；四个质量值互不相同时，过四点的三次多项式唯一存在，这是多项式插值的基本结论，可参考\BookChapterRef{chap:signal-analysis}{卷一的信号章}。取对数码率的理由是，原提案观察到在线性码率上积分会让高码率段主导结果，而同样比例的码率变化在对数尺度上等长；只在共同区间积分，是为了不外推任一条曲线。若B在整个区间上码率都是A的0.8倍，BD码率为$-20\%$，表示同等PSNR下平均节省两成码率；交换码率与质量的角色，可类似地得到同等码率下的平均PSNR差。视频编码比较BD码率时还要固定图像组长度、参考结构和测试序列，DVC的补充材料就按序列分别列出相对H.264的BD码率。拟合函数也会影响结果：原提案先讨论过有理函数，后来改用对数码率上的三次多项式，不同软件实现的拟合方式也不完全相同，所以BD码率应注明计算实现。

延迟衡量某一帧从拍到出结果要等多久。设输出第$t$帧需要读到第$t+r$帧，则
\begin{equation}
  L=\frac{r}{f}+T_{\mathrm{comp}}+T_{\mathrm{buf}}\text{，}
  \label{eq:vision-video-latency}
\end{equation}
三项依次是等待后帧、计算和显示或传输缓冲的时间。每秒30帧、等待两帧时，仅第一项就约为66.7毫秒。延迟与吞吐量不同：吞吐量描述每秒处理多少帧，批处理可以让吞吐很高，却让首帧等待很久。

\subsection{跨帧对齐误差带来的伪影}
\label{sec:vision-video-challenge-alignment}

多帧方法的收益来自把别的时刻对同一内容的观测搬到当前帧，搬错位置就会把错误的内容混进来。路牌上一个白色字母在第$t$帧位于横坐标100，在第$t-1$帧位于103；若估计的位移是2个像素，从邻帧取来的就是字母边缘与背景的混合，与当前帧合成后出现两条错开的边缘，即重影或双边缘。位移误差哪怕只有零点几个像素，在细笔画上也会变成可见的模糊；在压缩中，对应误差留给残差，块与块的运动不一致还会在块边界上形成方块感。

经典方法用显式的运动估计回答这个问题。块匹配为每个块搜索一个使匹配误差最小的位移（\ref{sec:vision-videofound-block-matching}节），计算规则、结果可解释，但一个块内只有一个位移，块里同时有车身和背景时必然有一部分对错；Horn--Schunck和Lucas--Kanade在亮度恒定和小运动假设下估计逐像素光流（\ref{sec:vision-videofound-horn-schunck}节、\ref{sec:vision-videofound-lucas-kanade}节），金字塔把大运动化为各层的小运动，可在运动边界、无纹理区域和亮度变化处仍会出错。经典去噪器VBM4D在沿运动轨迹堆叠图块时加入位置惩罚，并在相邻块差异过大时截断轨迹，用来限制错误对应进入滤波\scite{vision-vbm4d2012}。

学习方法从两处改变了这一环节。其一，光流网络按下游任务的损失端到端微调，学到的是对重建有用的对应，而不只是端点误差小的对应；其二，可变形对齐在初始对应附近读取多个位置，再由网络决定权重，使单个位置的误差不至于直接决定结果。代价是对应的含义变得不那么直观，错误也更难从中间结果上看出来。\ref{sec:vision-video-warping}节先讲反向取样和可见性检查，\ref{sec:vision-video-fusion}节讲RVRT怎样在光流附近多位置取样；编码中的对应误差由残差承担，见\ref{sec:vision-video-rate-allocation}节。

\subsection{遮挡与新显露区域带来的信息缺失}
\label{sec:vision-video-challenge-occlusion}

即使运动估计完全正确，邻帧也未必看到了当前帧需要的内容。一辆车从路牌前驶过，第$t$帧里刚被让出来的路牌左半部分，在第$t-1$帧中还被车身挡着；这块新显露的区域在过去帧中不存在可取的观测，任何从$t-1$搬来的值都是车身而不是路牌。插帧更极端：两侧都被遮挡的区域在任何输入中都没有记录。困难在于反向取样不会报错，它总能在邻帧中读出一个数值。

经典方法先检测，再退回。前后向一致性检查（\ref{sec:vision-video-warping}节式\eqref{eq:vision-video-fb-consistency}）把往返位移不能闭合的位置判为不可靠；VBM4D在遮挡或切镜导致相邻块差异超过阈值时中断轨迹，只用剩下的时间段滤波\scite{vision-vbm4d2012}；运动补偿时域滤波在补偿后的差异很大时减小时间权重，退回到当前帧（\ref{sec:vision-video-fusion}节）。在编码器中，新显露区域的块可以改用帧内预测，由率失真代价决定（\ref{sec:vision-video-rate-allocation}节）。

学习方法通常不显式输出遮挡掩码，而是让注意力或融合网络为不可靠的读取分配低权重，双向传播还可以从未来帧取得过去帧没有的内容。对于两侧都看不到的区域，网络只能依据训练数据中的先验补出内容；FILM的作者加入Gram矩阵损失，正是为了让大面积去遮挡区域的补全不至于过度模糊\scite{vision-film2022}。补出来的内容看起来合理，却不是观测，这一区别要在输出中保留。本章\ref{sec:vision-video-warping}节讲可见性检查，\ref{sec:vision-video-super-resolution}节讲双向传播怎样从未来帧取得内容，\ref{sec:vision-video-interpolation}节讲插帧中两侧都不可见区域的合成。

\subsection{时间一致性与逐帧质量的冲突}
\label{sec:vision-video-challenge-consistency}

逐帧做得最好的结果，连起来播放未必最好。若每帧独立去噪或放大，各帧补出的细节彼此独立，路牌上的纹理和字符边缘在相邻帧之间轻微跳动，播放时表现为闪烁；反过来，若强行让相邻输出相近，正常的运动会被拖慢，细节被平均掉，甚至留下前几帧的残影。前者在逐帧PSNR上可能更高，后者在扭曲误差上更低，二者不能同时由单一指标保证。

经典的运动补偿时域滤波用一个权重在两者之间折中：当前输出取当前观测与上一输出（经运动补偿）的加权和。\ref{sec:vision-video-fusion}节会推导，在对应正确、噪声独立的条件下，时间权重越大，噪声方差降得越多，但一次错误对应的残留也衰减得越慢。这个权重通常按补偿后的差异逐像素调整，本质上是在逐帧保真与时间平滑之间按局部可靠性分配。

学习方法把这种折中交给网络：循环传播让相邻输出共享同一个状态，输出自然更连贯；注意力权重相当于按内容自适应的时间权重。但网络的训练损失大多仍是逐帧重建误差，时间一致性只是间接得到；状态传播在带来连贯的同时也会传播错误。\ref{sec:vision-video-fusion}节讲经典滤波与RVRT的聚合和传播，\ref{sec:vision-video-temporal-quality}节说明怎样同时核验逐帧质量和时间连续性。

\subsection{等待范围与延迟约束}
\label{sec:vision-video-restoration-availability}

能读的帧越多，可用信息越多，但读未来的帧就要等待，读更多的帧就要缓存。输出第$t$帧时，逐帧处理只读当前帧；局部多帧处理读窗口$[t-l,t+r]$；循环传播把已处理的帧压缩进状态，单向时只含过去，双向时还含未来。按式\eqref{eq:vision-video-latency}，右侧等待$r$帧就至少多出$r/f$秒。离线方法若对整段做后向传播，右侧等待就是整段长度，应记录片段长度与缓存成本；把长视频切成有限片段处理，又会改变片段边界附近可用的帧。观察范围与信息到达条件的一般讨论见\ref{sec:vision-video-availability}节。

循环结构与在线因果是两回事。循环描述计算怎样复用：同一模块逐帧更新状态。因果描述输入何时可见：输出第$t$帧时是否读过$t$之后的帧。双向循环网络是循环的，却不是因果的；一个只读过去五帧的窗口网络不循环，却是因果的。把“循环”误读成“可在线”，是本章模型最常见的误用。

经典编码器早就面对同样的取舍。双向预测帧要参考显示顺序上更晚的帧，编码器必须先编码那一帧，于是解码顺序与显示顺序不同，解码端要多等若干帧才能输出；对延迟敏感的场合可以只用前向预测，DVC的单参考设置也属于这一类。学习式复原模型多为离线设计：BasicVSR和RVRT都要读到整段的末尾，FILM插中点必须等后一帧到达。\ref{sec:vision-video-neighbor-selection}节比较各种读取范围，\ref{sec:vision-video-codec-reference}节讲编码中的参考顺序，\ref{sec:vision-video-quality-rate-latency}节把质量与等待放在同一条曲线上比较。

\subsection{帧间冗余与码率分配}
\label{sec:vision-video-challenge-rate}

相邻帧的大部分内容相同，逐帧独立压缩会反复传输它们。帧间编码只传“怎样从已解码的帧预测当前帧”和“预测之后还差多少”两部分，但比特在两部分之间怎样分配并不显然：运动描述得越细，预测越准、残差越少，运动本身却越贵；新显露区域、切镜和曝光变化处，帧间预测可能还不如帧内预测。再加上首帧和随机访问点必须独立编码，一段视频的比特要在关键帧、运动、残差和模式选择之间分配。DVC的实验给出了一个具体量级：在其报告的三个工作点上，运动信息约占总码率的27\%到36\%，且质量要求越高，运动所占比例越低\scite{vision-dvc2019}。

经典编码标准用混合编码框架回答：块运动补偿负责帧间冗余，变换与量化负责空间冗余，熵编码负责统计冗余，编码器再用拉格朗日率失真代价为每个块选择预测模式和划分\scite{vision-rdo-video1998}。各模块按手工设计分别调优，标准只规定码流语法和解码过程。学习方法保留了这一结构，但把运动估计、运动压缩、运动补偿和残差变换换成网络，用一个率失真目标联合训练，使光流网络学会估计“便于压缩”的运动。\ref{sec:vision-video-codec}节按参考依赖、运动与残差编码、解码重建三步，先讲H.264/HEVC再讲DVC。

这五项挑战沿着数据流依次出现，本章也按这个顺序回答。前三项发生在“从邻帧取得信息”这一步：先决定读多远，再在对齐时处理位移误差与可见性，最后在聚合与传播时权衡时间平滑与逐帧细节，经典滤波与RVRT在这里对照。取得的信息随后要落到某个输出网格上，网格越远离原观测，遮挡和先验补全的问题越突出，于是从原网格恢复依次过渡到超分辨率和插帧。等待范围在每个模型中都要交代，在编码中则体现为参考顺序。码率分配是编码独有的问题，它要求对应能被传输和重现，所以放在复原之后。最后的核验把逐帧质量、时间一致性、码率和延迟放回同一协议下比较。

\section{跨帧有效信息的取得}
\label{sec:vision-video-information-fusion}

本节要解决的问题是：给定低质量序列$\bm{Y}_1,\dots,\bm{Y}_n$和目标时刻$t$，为第$t$帧的每个位置取得其他时刻对同一内容的可靠观测，输出一份与第$t$帧网格对齐的特征或画面。难点正是上一节的前三项挑战：同一内容在不同帧的位置不同，需要对应，对应有误差就产生伪影；对应正确时内容仍可能被遮挡或曝光改变，需要判断可用性；为了看得更远而传递的状态，又会在平滑与细节之间取舍，并积累错误。

这些难点按处理顺序依次出现，所以本节按步骤展开：先决定读哪些帧，再把它们对齐到目标网格，最后聚合并向后传播。对同一难点，不同方法给出不同回答。读取范围上，滑动窗口在每个输出时刻重新读取固定邻帧，范围有限但互不干扰；循环方法把历史压缩进状态，看得远但串行且会传播错误；RVRT在片段内并行联合处理几帧，片段之间再循环传递，是两者的折中。对齐上，块匹配为每个块取一个位移，光流为每个像素取一个位移，可变形对齐在光流附近取多个位置再加权；越往后，对单个位移误差的容忍越强，需要学习的量也越多。聚合上，经典的运动补偿时域滤波用一个按局部差异调整的权重递归平均，VBM4D沿运动轨迹把相似图块堆成组再协同收缩，RVRT则用注意力在相邻片段之间读取并传播特征。RVRT放在最后一步完整讲解，因为它的主要改动落在对齐后怎样聚合与传播；以双向循环为核心的BasicVSR放在\ref{sec:vision-video-basicvsr}节，与超分辨率任务一起讲。

\subsection{邻帧范围的选择}
\label{sec:vision-video-neighbor-selection}

选帧要回答：哪些时刻可能观察到第$t$帧缺少的内容，读取它们要付出多少等待和计算。邻帧的价值不在“多一张图”，而在它是否看到了目标位置需要的内容。

固定窗口直接读取当前帧前后各$k$帧，计算范围明确，各输出时刻相互独立，可以并行；代价是每帧特征被相邻$2k+1$个输出重复读取，且窗口外的清晰观测无法利用。单向循环让历史信息沿时间累积，同一模块复用于各帧，每帧只处理一次；但状态容量有限，会压缩细节并传播错误，序列开头的帧也几乎得不到上下文。双向循环再加一条从后向前的状态，使每帧同时获得整段的前后信息，代价是等待整段输入并保存中间状态。片段循环先把相邻$N$帧组成片段并行处理，再在片段之间传递状态，介于两者之间。表\ref{tab:vision-video-range-strategies}按同一组维度比较这四种范围。

\begin{table}[htbp]
  \centering
  \begin{tabularx}{\linewidth}{@{}lXXXX@{}}
    \toprule
    范围策略 & 可见范围 & 右侧等待 & 每帧特征计算次数 & 主要风险\\
    \midrule
    固定窗口 & $[t-k,t+k]$ & $k$帧 & 约$2k+1$次 & 错过窗口外的清晰观测\\
    单向循环 & $[1,t]$ & 0帧 & 1次 & 状态遗忘与错误传播；开头帧上下文少\\
    双向循环 & $[1,n]$ & 至整段末尾 & 每个方向1次 & 等待整段；缓存全部中间状态\\
    片段循环 & 片段内联合，片段间传播 & 依传播方向而定 & 每层1次 & 片段内运动过大时联合处理失效\\
    \bottomrule
  \end{tabularx}
  \caption{邻帧范围的四种策略。“每帧特征计算次数”指一帧的特征在整段处理中被重新提取的次数，固定窗口可以缓存特征来减少重复，但会增加内存。片段循环若交替方向，也会读到未来帧。}
  \label{tab:vision-video-range-strategies}
\end{table}

用路牌序列看范围的影响。某块车牌在$t-3$帧清晰、$t$帧被雨滴遮住、$t+2$帧再次出现，半径$k=1$的窗口看不到两次清晰观测；单向循环可能把$t-3$帧的特征带来，却不能保证细笔画在三次更新后仍保存在状态中；双向循环能同时利用$t-3$与$t+2$，但必须等到$t+2$帧到达。若中间发生切镜，继续传播旧场景特征反而有害，系统可按\ref{sec:vision-videofound-shot-boundary}节的镜头边界检测重置状态。BasicVSR的作者把测试序列切成更短片段处理，观察到可用时间范围缩短时重建质量下降，据此认为远处帧的信息不应忽略\scite{vision-basicvsr2021}。这说明范围有收益，但收益仍要和遮挡、切镜与等待一起检查。

\subsection{邻帧内容的对齐}
\label{sec:vision-video-warping}

选定邻帧$s$后，要回答“目标帧位置$\bm{p}$的内容在邻帧的哪里”。输入是第$t$帧与第$s$帧（或它们的特征），输出是把邻帧内容搬到第$t$帧网格后的结果。难点是位移不是整数，估计可能有误，被搬来的位置也可能根本不可见。对齐有三条由粗到细的路线：按块取一个位移，按像素沿光流读取单个位置，或在光流附近学习多个取样位置再加权；越往后越能容忍局部误差，但需要估计或学习的量也越多。

最早被大量使用的是块匹配。它把当前帧切成固定大小的块，在邻帧的搜索窗口内找使绝对差之和最小的位置，搜索方法与代价见\ref{sec:vision-videofound-block-matching}节。块匹配的结果是一组块位移，同一块内所有像素共用一个位移，这正是视频编码标准采用的运动表示（\ref{sec:vision-video-rate-allocation}节）。它在块内运动一致时有效，块内同时包含前景和背景、或发生旋转和缩放时，总有一部分像素对错。逐像素的光流消除了块的限制：Horn--Schunck用全局平滑项、Lucas--Kanade用局部窗口最小二乘，在亮度恒定和小位移的假设下求每个像素的位移，配合图像金字塔处理大运动（\ref{sec:vision-videofound-horn-schunck}节、\ref{sec:vision-videofound-lucas-kanade}节）。

本章三个学习模型RVRT、BasicVSR和DVC都用同一个光流网络SPyNet初始化对应，它把Lucas--Kanade式的由粗到细金字塔直接搬进了网络\scite{vision-spynet2017}。SPyNet先把两帧$\bm{I}_1,\bm{I}_2$各建成5级图像金字塔，第$k=0$级最粗，每级分辨率是下一级的一半。在第$k$级，它把上一级的光流放大两倍得到$u(\bm{V}_{k-1})$，用它把第二帧反向取样到第一帧的网格，再由一个小卷积网络$G_k$估计剩余的位移：
\begin{equation}
  \bm{v}_k=G_k\bigl(\bm{I}^{k}_1,\ \mathcal{W}(\bm{I}^{k}_2,u(\bm{V}_{k-1})),\ u(\bm{V}_{k-1})\bigr),\qquad
  \bm{V}_k=u(\bm{V}_{k-1})+\bm{v}_k\text{，}
  \label{eq:vision-video-spynet}
\end{equation}
最粗一级的初始光流取零。$G_k$的输入是两幅RGB图像和两通道光流拼成的8通道张量，由5层$7\times7$卷积组成，输出通道依次为32、64、32、16、2，除最后一层外每层后接ReLU。原文用Flying Chairs等合成数据上的端点误差逐级训练这5个网络，每级以上一级的网络初始化。由于每一级只需估计很小的修正（原文称通常不到一个像素），$G_k$可以很小；这与经典金字塔光流“粗层看大运动、细层补小运动”的分工相同，区别在于每一级的更新规则由数据学到，而不是由亮度恒定方程推出。

无论位移来自块匹配、经典光流还是网络，取值方式都一样。设目标帧位置为$\bm{p}=(x,y)^\top$，由目标帧$t$指向邻帧$s$的光流为$\bm{u}_{t\to s}(\bm{p})$，它可以由SPyNet、\ref{sec:vision-video-raft}节的RAFT或其他方法估计。为了在目标网格中重建同一内容，应到邻帧的位置$\bm{p}+\bm{u}_{t\to s}(\bm{p})$取值。方向约定很重要：若误用$s\to t$的位移而不转换，通常不能靠简单取负得到正确采样位置，因为两者定义在不同帧的网格上。

\term{反向取样}（backward warping）逐个遍历输出位置，到源图像查询内容。与之相对的前向搬运把源像素推到目标位置，会形成空洞或多对一写入冲突；反向取样保证每个输出位置恰好得到一个值。查询位置通常不是整数。若横纵坐标分别为$j+\alpha,i+\beta$，其中$0\leq\alpha,\beta\leq1$，双线性插值用周围四个像素计算某个通道的值：
\begin{equation}
  \begin{aligned}
    \widetilde y={}&(1-\alpha)(1-\beta)y_{i,j}
       +\alpha(1-\beta)y_{i,j+1}\\
       &+(1-\alpha)\beta y_{i+1,j}
       +\alpha\beta y_{i+1,j+1}\text{。}
  \end{aligned}
  \label{eq:vision-video-bilinear}
\end{equation}
这里的四个$y$都来自邻帧$s$同一通道，省略了时间与通道标记。插值核的一般讨论见\ref{sec:vision-formation-interpolation}节。权重随距离连续变化且和为1，因此结果对位移可导，光流网络可以通过取样结果接收梯度；但插值只说明怎样读数，不保证读到的位置确实可见。

在教学网格中，四角值从左上到右下为$0.2,0.6,0.4,1.0$，取样位置为$\alpha=0.25,\beta=0.5$。对应权重为$0.375,0.125,0.375,0.125$，得到
\[
  \widetilde y=0.375(0.2)+0.125(0.6)
       +0.375(0.4)+0.125(1.0)=0.425\text{。}
\]
图\ref{fig:vision-bilinear-warp}把位置与权重放在同一网格中。若移动字母的位移估计少了一个像素，取来的可能是字母边缘与背景的混合，和当前帧融合后就会产生双边缘。

\begin{figure}[htbp]
  \centering
  % 四角与查询点均为正文可复算的教学数值。
\begin{tikzpicture}[x=1cm,y=1cm,font=\figurefont\small]
  \draw[BookRule,line width=0.6pt] (1,0) rectangle (4.5,-2.8);
  \foreach \x/\y/\val/\wt/\anchor in {
    1/0/0.2/0.375/south west,
    4.5/0/0.6/0.125/south east,
    1/-2.8/0.4/0.375/north west,
    4.5/-2.8/1.0/0.125/north east}{
    \fill[BookBlue] (\x,\y) circle (2pt);
    \node[anchor=\anchor] at (\x,\y) {\val\ (\wt)};
  }
  \draw[BookMuted,dashed,line width=0.6pt] (1.875,0) -- (1.875,-2.8);
  \draw[BookMuted,dashed,line width=0.6pt] (1,-1.4) -- (4.5,-1.4);
  \fill[BookTeal] (1.875,-1.4) circle (2.5pt);
  \node[above right] at (1.875,-1.4) {取样点};
  \draw[-{Stealth[length=2mm]},BookInk,line width=0.8pt] (0.6,0.65) -- (4.8,0.65);
  \node[anchor=west] at (4.8,0.65) {$x$};
  \draw[-{Stealth[length=2mm]},BookInk,line width=0.8pt] (0.5,0.45) -- (0.5,-3.2);
  \node[below] at (0.5,-3.2) {$y$};
  \node[anchor=west,align=left] at (5.7,-0.6) {$\alpha=0.25$\\$\beta=0.5$};
  \node[anchor=west,align=left] at (5.7,-2.0) {四角加权\\$\widetilde y=0.425$};
\end{tikzpicture}

  \caption{双线性取样的教学网格。横坐标向右、纵坐标向下；查询点位于两行的中间、偏左四分之一处。四角旁括号内是插值权重，权重和为1，输出为0.425。图中线段表示插值邻域，不表示物体真实运动轨迹。}
  \label{fig:vision-bilinear-warp}
\end{figure}

图像对齐移动像素值，特征对齐移动网络提取的表示。特征对齐的好处是取样误差造成的模糊可以由后续网络部分补偿；BasicVSR报告，把特征对齐换成图像对齐会使PSNR下降0.17 dB\scite{vision-basicvsr2021}。两者都要统一坐标和尺度：在下采样$2^\ell$倍的特征上取样时，位移数值也要除以$2^\ell$。越界位置需有明确的填充及有效性标记。

可见性要单独判断。邻帧中的位置可能已被前景挡住，插值照样给出一个数值。常用的筛选线索是前后向一致性：从$\bm{p}$沿$\bm{u}_{t\to s}$到达邻帧，再沿邻帧到目标帧的光流$\bm{u}_{s\to t}$回来，若内容可见，应大致回到出发点，即
\begin{equation}
  \bigl\lVert\bm{u}_{t\to s}(\bm{p})
   +\bm{u}_{s\to t}\bigl(\bm{p}+\bm{u}_{t\to s}(\bm{p})\bigr)\bigr\rVert_2
  \leq\delta\text{，}
  \label{eq:vision-video-fb-consistency}
\end{equation}
其中$\delta$是以像素计的容许阈值，$\bm{u}_{s\to t}$在非整数位置同样用双线性插值读取。教学数值：若$\bm{u}_{t\to s}(\bm{p})=(3,0)^\top$，回程读到$(-2.6,0.1)^\top$，和为$(0.4,0.1)^\top$，长度约0.41像素；取$\delta=1$时判为一致。若该位置在邻帧被遮挡，回程光流描述的是遮挡物的运动，和可能达到数个像素，从而被排除。这个检查依赖两份光流本身，两者同时出错时仍会漏判。

即使对应和可见性都成立，曝光差异和反光也可能改变亮度，光度校正或可靠性权重属于额外判断。单一位置取样对光流误差十分敏感，于是出现了第三条路线：\term{可变形对齐}（deformable alignment）在初始对应附近由网络预测若干偏移，到多个位置取样，再加权合成，让网络自己在附近寻找正确内容。它的代价是偏移由谁预测、权重怎样形成都必须交代清楚；RVRT的引导可变形注意力是这一路线的具体实现，见下一小节。

\subsection{对齐信息的聚合与传播}
\label{sec:vision-video-fusion}
\label{sec:vision-video-rvrt}

对齐之后，要把当前帧与多份对齐结果合成一份输出，并决定把什么传给后续帧。输入是当前帧（或其特征）和若干对齐后的邻帧，输出是更新后的当前结果与传递状态。难点在于，错误对应即使只占少数，也可能主导合成结果；而传递出去的错误会在后续每一帧重复出现。

先用一个教学例子说明为何不能简单平均。当前观测值为0.40，对齐后的邻帧值为0.425，假设可靠性分别为1和3，加权合成为$(0.40+3\times0.425)/4=0.41875$。若错误对应带来0.90且仍赋予权重3，结果就变成0.775。分母归一化不能消除错配，多帧汇总是否有益取决于对应及权重是否可靠。下面三种方法都在回答“权重从哪里来”：运动补偿时域滤波用补偿后的差异逐像素决定，VBM4D用图块之间的相似度决定，RVRT用学到的注意力决定。

沿运动轨迹做时间滤波是经典视频降噪的一大类方法\scite{vision-brailean1995}。最简单的形式是递归滤波：把上一帧的输出按运动补偿到当前网格，记为$\widetilde{\bm{X}}_t=\mathcal{W}(\widehat{\bm{X}}_{t-1},\bm{u}_{t\to t-1})$，再与当前观测加权平均，
\begin{equation}
  \widehat{\bm{X}}_t(\bm{p})=\bigl(1-a_t(\bm{p})\bigr)\bm{Y}_t(\bm{p})
    +a_t(\bm{p})\,\widetilde{\bm{X}}_t(\bm{p})\text{，}
  \label{eq:vision-video-mctf}
\end{equation}
其中$a_t(\bm{p})\in[0,1)$是时间权重。递归的好处是只需保存一帧输出，却相当于对整段历史做了指数加权平均。

时间权重怎样影响降噪，可以在理想条件下算出来。设场景静止或运动补偿完全正确，真值在该位置为常数$x$，各帧噪声独立、均值为零、方差为$\sigma^2$，权重取常数$a$。记误差$e_t=\widehat{X}_t(\bm{p})-x$，由式\eqref{eq:vision-video-mctf}得$e_t=(1-a)\varepsilon_t+a\,e_{t-1}$，其中$\varepsilon_t$是第$t$帧的噪声。$e_{t-1}$只依赖第$t$帧之前的噪声，与$\varepsilon_t$独立，所以方差相加：
\[
  v_t=(1-a)^2\sigma^2+a^2v_{t-1}\text{。}
\]
当$|a|<1$时，这个线性递推收敛到不动点$v=(1-a)^2\sigma^2+a^2v$，解得
\begin{equation}
  v=\frac{(1-a)^2}{1-a^2}\sigma^2=\frac{1-a}{1+a}\sigma^2\text{。}
  \label{eq:vision-video-mctf-variance}
\end{equation}
第二个等号用了$1-a^2=(1-a)(1+a)$。$a=0.5$时方差降到$\sigma^2/3$，$a=0.8$时降到$\sigma^2/9$，与直接平均9帧相当。代价体现在错误的残留上：若某一帧因对应错误混入了偏差$\delta$，它在$k$帧之后仍以$a^k\delta$的大小留在输出里。$a=0.8$时，三帧后还剩$0.512\delta$；$a=0.5$时只剩$0.125\delta$。时间权重越大，降噪越强，残影也越久，这正是上一节第三项挑战的定量形式。

因此实际的滤波器让权重随局部差异变化：补偿后的上一输出与当前观测差得越多，越可能是对应错误或遮挡，权重就越小。算法\ref{alg:vision-video-mctf}给出一个教学版本，权重函数$a=a_{\max}\max(0,1-d/\kappa)$是为说明机制而选的简单形式，不对应某一篇论文的具体设计。

\begin{algorithm}[运动补偿时域递归滤波（教学版本）]
  \label{alg:vision-video-mctf}
  \begin{algorithmic}[1]
    \Require 观测序列$\bm{Y}_1,\dots,\bm{Y}_n$；最大时间权重$a_{\max}\in[0,1)$；差异尺度$\kappa>0$；光流估计方法
    \Ensure 滤波结果$\widehat{\bm{X}}_1,\dots,\widehat{\bm{X}}_n$
    \State $\widehat{\bm{X}}_1\gets\bm{Y}_1$
    \For{$t=2,\dots,n$}
      \State 估计光流$\bm{u}_{t\to t-1}$，计算$\widetilde{\bm{X}}_t\gets\mathcal{W}(\widehat{\bm{X}}_{t-1},\bm{u}_{t\to t-1})$
      \For{每个位置$\bm{p}$}
        \State $d\gets\lvert\bm{Y}_t(\bm{p})-\widetilde{\bm{X}}_t(\bm{p})\rvert$
        \State $a\gets a_{\max}\max(0,1-d/\kappa)$
        \State $\widehat{\bm{X}}_t(\bm{p})\gets(1-a)\bm{Y}_t(\bm{p})+a\,\widetilde{\bm{X}}_t(\bm{p})$
      \EndFor
    \EndFor
  \end{algorithmic}
\end{algorithm}

用一个像素复算。取$a_{\max}=0.8$、$\kappa=0.2$，路牌上某个已对准的像素真值约为0.5，观测依次为0.50、0.42、0.55。第2帧差异$d=0.08$，权重$a=0.8\times0.6=0.48$，输出$0.52\times0.42+0.48\times0.50=0.4584$；第3帧$d=|0.55-0.4584|=0.0916$，$a=0.8\times(1-0.458)=0.4336$，输出$0.5664\times0.55+0.4336\times0.4584\approx0.5103$。两次输出都比观测更接近0.5。第4帧一辆车驶入该位置，观测跳到0.90，$d=0.39>\kappa$，权重变为0，输出直接取0.90，不会把路牌的值拖成残影。反过来，若这次跳变其实是一个脉冲噪声，滤波器同样会把它原样保留：固定阈值区分不了“新内容”和“大噪声”。

VBM4D把非局部均值（\ref{sec:vision-transforms-nonlinear-denoising}节）“在别处找相似图块”的思路推广到视频\scite{vision-vbm4d2012}。它先沿运动轨迹逐帧跟踪一个图块，得到一个三维时空体；再在附近找与之相似的其他时空体，沿第四个维度堆叠成“组”。组内有三种相关：块内的局部相关、沿轨迹的时间相关、不同体之间的非局部相似。对组做可分离的四维去相关变换，收缩变换系数，再逆变换，把每个体的估计放回原位置并加权聚合，这一步称为协同滤波。整个过程分两级：第一级用硬阈值收缩得到初步估计，第二级以初步估计为引导做经验维纳滤波。轨迹在相邻块差异超过阈值时中断，以应对遮挡和切镜。同一框架还用于去块效应和增强。

两种经典方法的设计思想都在学习方法中延续了下来。按局部差异调节的时间权重，对应注意力中由查询与键的相似度决定的权重；把多个相似候选放在一起再加权，对应在多个取样位置上的注意力聚合；轨迹中断，对应注意力为不可靠的候选分配低权重。差别在于，经典方法的相似度量、阈值和变换是手工设定的，学习方法从配对数据中学到这些量，可以处理更复杂的退化，代价是需要配对数据、计算量大得多，而且行为更难逐项解释。

学习方法在组织聚合与传播时也有两种极端。并行方法把整段视频的所有帧一起送入网络，用注意力在帧间交换信息；它的帧间融合充分，但模型大、对长序列内存消耗巨大。循环方法逐帧更新共享参数的状态，模型小；但它每一步只能把一帧信息压进状态，RVRT的作者指出这种串行处理会造成信息损失和噪声放大\scite{vision-rvrt2022}。\ref{sec:vision-video-basicvsr}节的BasicVSR属于后一种。

RVRT（recurrent video restoration Transformer）的主要想法是：把视频切成每段$N$帧的相邻片段，片段内部的几帧并行联合更新，片段之间按循环方式传递特征；跨片段对齐使用引导可变形注意力，在光流给出的位置附近为每帧读取多个候选。$N=1$时它退化为逐帧循环，$N=n$时它变成整段并行，因此可以在两种极端之间选择折中。片段内联合处理让相邻帧直接交换信息，缓解了逐帧压缩造成的损失；片段间循环又避免了整段并行的内存开销。

RVRT的主干由Swin Transformer层搭成。一个Swin层把特征划分成不重叠的$M\times M$窗口，在每个窗口内做多头自注意力，相邻层交替使用规则划分和平移半个窗口的划分，使信息能跨窗口流动；窗口注意力的一般计算可参考\BookChapterRef{chap:transformer}{卷五的Transformer章}。设窗口内的词元矩阵为$\bm{Z}$，一层的计算顺序是
\begin{equation}
  \bm{Z}'=\bm{Z}+\operatorname{MSA}\bigl(\operatorname{LN}(\bm{Z})\bigr),\qquad
  \bm{Z}''=\bm{Z}'+\operatorname{MLP}\bigl(\operatorname{LN}(\bm{Z}')\bigr)\text{，}
  \label{eq:vision-video-swin-layer}
\end{equation}
其中$\operatorname{LN}$是层归一化，$\operatorname{MSA}$是窗口内的多头自注意力（带可学习的相对位置偏置），$\operatorname{MLP}$是中间夹GELU的两层全连接网络。若干个Swin层串联，末端接一层卷积，再与块的输入相加，就是\term{残差Swin Transformer块}（residual Swin Transformer block, RSTB），它来自图像复原模型SwinIR\scite{vision-swinir2021}：$\bm{F}_{\mathrm{out}}=\operatorname{Conv}\bigl(\mathrm{STL}_L\circ\cdots\circ\mathrm{STL}_1(\bm{F}_{\mathrm{in}})\bigr)+\bm{F}_{\mathrm{in}}$，$\mathrm{STL}$表示一个Swin层。卷积为注意力补充平移等变的局部处理，残差连接让不同深度的特征直接通到后面。

RVRT的输入是$n$帧低质量视频，形状为$n\times H\times W\times3$，整体分三段。浅层特征提取先用一层卷积，再接1个含2个Swin层的RSTB，得到$n\times H\times W\times C$的特征；超分辨率配置取$C=144$，去模糊和去噪配置取$C=192$。去模糊、去噪配置的输出与输入同尺寸，原实现另加两层步长卷积先降低特征的空间尺寸以控制计算，最后再恢复。随后特征按$N$帧一组切成片段，原文取$N=2$。中段是$L=4$个循环特征细化模块，最后经1个含2个Swin层的RSTB整理特征，由像素重排层（见\ref{sec:vision-video-basicvsr}节）把通道转换为$s$倍空间尺寸，输出$n\times sH\times sW\times3$的视频。表\ref{tab:vision-video-rvrt-config}列出各部分的配置。

\begin{table*}[tbp]
  \centering
  \small
  \begin{tabularx}{\linewidth}{@{}lXrrrrl@{}}
    \toprule
    模块 & 窗口或取样范围 & 每窗口词元数 & 层数 & 宽度 & 头数 & MLP倍率\\
    \midrule
    浅层特征提取 & 1个RSTB，$8\times8$空间窗口 & 64 & 2 & $C$ & 6 & 论文未给出\\
    循环特征细化（每个模块） & 1层融合卷积，2个修改的RSTB，$N\times8\times8$时空窗口 & 128 & 4 & $C$ & 6 & 论文未给出\\
    引导可变形注意力（每个模块） & 前一片段2帧，每帧9个候选位置 & 18 & 1 & 查询投影到$2C$ & 12组，12头 & 论文未给出\\
    重建 & 1个RSTB，$8\times8$空间窗口，像素重排 & 64 & 2 & $C$ & 6 & 论文未给出\\
    \midrule
    全模型 & 4个细化模块，合计16个修改的Swin层 & — & 20 & 144或192 & — & —\\
    \bottomrule
  \end{tabularx}
  \caption{RVRT的配置，取自原论文的实验设置，片段长度$N=2$。宽度$C$在超分辨率配置中为144，在去模糊和去噪配置中为192。“层数”指Swin层的个数，全模型一行把浅层、细化和重建三段相加。原论文报告的参数量为：超分辨率1080万（低质量输入$320\times180$时运行时间183毫秒），去模糊1360万，去噪1280万。Swin层和引导可变形注意力中多层感知机的隐藏层倍率，论文没有给出。}
  \label{tab:vision-video-rvrt-config}
\end{table*}

核心是4层循环特征细化。每层按时间顺序依次处理各片段；第$k$个片段在第$i$层的计算分三步。第一步，用引导可变形注意力把本层已经算好的前一片段特征对齐到当前片段。第二步，把对齐结果与当前片段在第$0$到$i-1$层的特征沿通道拼接，用一层卷积融合回$C$个通道。第三步，用2个修改过的RSTB细化，每个含2个Swin层：它们把二维$8\times8$窗口扩展为$N\times8\times8$的时空窗口，使片段内各帧在同一窗口中互相注意；$N=2$时每个窗口含128个位置，每层6个头。偶数层把视频序列反转后再处理，于是各层交替从前向后、从后向前传播。

\term{引导可变形注意力}（guided deformable attention, GDA）负责第一步的对齐。它要在前一片段的$N$帧中，为当前帧位置$\bm{p}$找到相关内容。SPyNet先由低质量输入估计光流$\bm{u}_{t\to s}$，用它把支持帧特征预对齐；一个由若干卷积层和ReLU组成的小网络读取当前特征、预对齐的支持特征和光流三者的拼接，为每个支持帧$s$预测$M$个偏移$\Delta\bm{p}_{s,m}$，原实现取$M=9$。第$m$个取样位置为
\[
  \bm{p}_{s,m}=\bm{p}+\bm{u}_{t\to s}(\bm{p})
                   +\Delta\bm{p}_{s,m}\text{，}
\]
非整数处沿用式\eqref{eq:vision-video-bilinear}取样。光流把搜索范围缩小到正确位置附近，偏移在附近微调取样点，注意力再决定每个取样值占多大比重，三者分工不同。

令当前查询为$\bm{q}$，第$(s,m)$个取样位置上的键和值为$\bm{k}_{s,m}$、$\bm{v}_{s,m}$，单个注意力头计算
\begin{equation}
  \begin{aligned}
    a_{s,m}&=
      \frac{\exp(\langle\bm{q},\bm{k}_{s,m}\rangle/\sqrt d)}
      {\sum_{s',m'}\exp(\langle\bm{q},\bm{k}_{s',m'}\rangle/\sqrt d)},\\
    \bm{r}&=\sum_{s,m}a_{s,m}\bm{v}_{s,m}\text{，}
  \end{aligned}
  \label{eq:vision-video-gda}
\end{equation}
其中$d$是查询和键的通道维度，求和遍历前一片段的$N$帧和每帧$M$个位置，$N=2$、$M=9$时共18个候选。查询由当前片段在上一层的特征投影得到，键从前一片段在上一层的特征中取样，值从前一片段在本层已更新的特征中取样；比较依据来自同一层次的特征，待传递的内容则是最新结果。原实现先投影再取样，以减少重复计算，并把查询投影到$2C$维，作者称这能略微提升效果。通道被分成12个可变形组分别取样，并用12个注意力头分别计算，所有组与头的结果拼接后，再经带残差连接的两层感知机交换通道信息，因为式\eqref{eq:vision-video-gda}只在空间位置之间加权。每层结束时，光流加上$M$个偏移的平均值，作为下一层的引导，使引导位置逐层修正。

回到车牌例子。光流给出的初始位置可能偏向笔画一侧，多个取样点可以同时覆盖笔画两侧，再按与当前特征的相似度聚合。若所有取样点都落在雨滴上，注意力也没有真实字符可读，它只能在给定候选中分配权重。交替方向的多层传播还意味着第$t$帧的输出读过$t$之后的帧，RVRT虽然有循环结构，仍属于离线方法。

RVRT用低质量序列与高质量目标配对训练，对所有输出帧施加逐元素的Charbonnier损失：
\begin{equation}
  L_{\mathrm{rec}}=\frac{1}{|\Omega|}\sum_{(t,i,j,k)\in\Omega}
    \sqrt{\bigl(\widehat{x}_{t,i,j,k}-x_{t,i,j,k}\bigr)^2+\epsilon^2}\text{，}
  \label{eq:vision-video-charbonnier}
\end{equation}
其中$\Omega$是所有帧、空间位置和通道的下标集合，$\epsilon>0$是小常数，RVRT取$10^{-3}$。误差远大于$\epsilon$时，每项约等于绝对误差，对少数大误差不像平方误差那样过度敏感；误差接近零时，函数仍光滑可导。原实现从$256\times256$的高质量块随机裁剪，在REDS上每个样本取30帧，在Vimeo-90K上取14帧，用Adam以批量8训练60万步，学习率从$4\times10^{-4}$按余弦退火下降。SPyNet以预训练权重初始化，前3万步冻结，之后学习率降为其余参数的四分之一；冻结是为了避免随机初始化的其他模块在训练初期把光流破坏。去噪模型按非盲设定训练，噪声标准差在0到50之间均匀抽取。推理时整段输入，各层逐片段运行，不针对每段视频重新拟合。

原文的消融实验在缩减一半修改RSTB、用12帧训练的简化模型上进行，可以看出各设计的作用：片段长度从1增至2，REDS上的PSNR从31.98升到32.10 dB，增至3时为32.07 dB，作者将不再提升归因于片段内运动较大和光流不准；去掉引导可变形注意力中的光流引导，PSNR降到30.99 dB。原论文在超分辨率、去模糊和去噪上验证了这一设计，这些结果不直接证明低照度增强效果。对切镜或错误传播，应分别查看帧内清晰度与后续连续性；重复传递的错误特征来自同一错误，不会因为经过多个片段就变成多份独立证据。

表\ref{tab:vision-video-aggregation-methods}在同一去噪任务下比较三种聚合方式。从运动补偿时域滤波到VBM4D，改变的是候选的组织方式：从“上一帧同一位置的一个值”扩大到“多条轨迹上的多个相似图块”，换来更强的降噪，代价是块匹配与变换的计算量。从VBM4D到RVRT，改变的是权重的来源：从手工的相似度和阈值换成学到的注意力，换来对复杂退化的适应能力，代价是需要配对训练数据和大得多的模型，以及输出依赖于训练分布。

\begin{table}[htbp]
  \centering
  \small
  \begin{tabularx}{\linewidth}{@{}lXXX@{}}
    \toprule
    比较维度 & 运动补偿时域滤波 & VBM4D & RVRT\\
    \midrule
    对应来源 & 每帧一份光流或块位移 & 带位置惩罚的块轨迹跟踪 & SPyNet光流加学到的多位置偏移\\
    聚合权重 & 按补偿后差异设定的时间权重 & 相似图块分组后的变换域收缩 & 查询与键相似度的注意力权重\\
    可见性处理 & 差异过大时权重降为0 & 差异超过阈值时截断轨迹 & 隐式地为不可靠候选分配低权重\\
    时间范围 & 递归，等效于指数加权的全部过去 & 轨迹覆盖的有限前后帧 & 交替方向传播，覆盖整段\\
    需要训练 & 否 & 否 & 是，需要配对数据\\
    \bottomrule
  \end{tabularx}
  \caption{三种跨帧聚合方式的比较。比较的是机制，不是在某一数据集上的质量排序；运动补偿时域滤波的设定为算法\ref{alg:vision-video-mctf}的教学版本。}
  \label{tab:vision-video-aggregation-methods}
\end{table}

\section{目标画面的恢复与采样扩展}
\label{sec:vision-video-restoration-output}

取得跨帧特征后，要把它转换成目标网格上的画面。输入是对齐聚合后的特征，或两帧原图；输出由\ref{sec:vision-video-degradation}节规定的网格决定。难点随网格而变：原网格上每个位置都有退化观测，主要难在区分细节与噪声；更高分辨率网格上新增的位置只有邻帧的间接观测；中间时刻则没有任何直接观测，只能根据两侧推算运动后再合成，遮挡带来的信息缺失在这里最突出。

这三类输出由同一组模块的不同组合实现。原网格恢复在聚合特征后接重建头即可，以RVRT的去噪、去模糊配置为例，并与VBM4D这类经典多帧滤波对照；更高分辨率恢复还需要上采样，并且更依赖长范围的亚像素观测，BasicVSR用双向循环传播和像素重排回答这一点；中间时刻恢复则把对齐方向反过来，从待生成的时刻出发去两侧读取，FILM在多尺度上学习这种读取位移，不需要光流标注。三者都使用反向取样，差别在于取样的起点、位移的来源以及输出位置有无直接观测。

\subsection{原采样画面的恢复}
\label{sec:vision-video-native-restoration}

原采样恢复要在帧时刻和空间尺寸都不变的条件下去除噪声、模糊或改善亮度。邻帧只提供补充证据，当前帧本身仍是主要观测。图像去噪的退化模型、条件均值解释和残差学习已在\ref{sec:vision-restoration-denoising}节讲过；视频版本的新增问题在于训练配对和时间稳定。

训练样本应说明目标是怎样取得的，例如对清晰序列合成规定噪声，或使用已对齐的模糊与清晰拍摄配对。若目标帧与输入帧错位一两个像素，网络会受到互相冲突的监督，学会的往往是模糊折中；合成退化与真实夜间成像不一致，也会影响迁移。

多帧去噪的收益来自邻帧对同一内容的独立噪声观测：在对应正确、噪声独立的理想条件下，平均$K$个观测可以把噪声方差降到$1/K$，式\eqref{eq:vision-video-mctf-variance}的递归滤波是同一原理的指数加权版本；对应错误或噪声在帧间相关时，这一收益随之减少。经典方法中，VBM4D不需要训练，只要给定噪声标准差就能运行，并且同一框架可以去除压缩造成的块效应\scite{vision-vbm4d2012}；它的局限是退化模型固定为加性噪声或块效应，真实夜间录像中随亮度变化的噪声、运动模糊和曝光跳变不在其模型之内。

RVRT的去噪和去模糊配置沿用\ref{sec:vision-video-rvrt}节的跨片段机制，宽度取$C=192$，浅层用两层步长卷积降低特征尺寸，重建时再恢复到原尺寸；输出尺度为1，并不意味着中间特征都保持原尺寸。原文报告，去噪模型约有1280万个参数，在$1280\times720$的输入上运行约0.2秒，去模糊模型约有1360万个参数、约0.3秒\scite{vision-rvrt2022}。去噪模型按非盲设定训练，即把噪声标准差作为已知量提供给模型；这与VBM4D需要噪声水平的条件相同，真实场景中仍要先估计它。

夜间路牌提亮后，要检查的不只是平均亮度，还包括相邻帧是否明暗跳变、白色字符是否偏色、运动边缘是否拖尾。降噪抹掉细笔画、锐化改变字符、局部增强放大闪烁，都是网格不变时仍可能出现的失真。增强服务观看时可以接受的改变，用于文字读取或亮度测量时可能不可接受。

\subsection{更高分辨率画面的恢复}
\label{sec:vision-video-super-resolution}
\label{sec:vision-video-basicvsr}

视频超分辨率的输入是$n$帧$h\times w$低分辨率画面，输出是$n$帧$sh\times sw$画面。难点在于新增像素没有直接观测，单帧放大只能依靠先验补出细节；多帧方法则希望从邻帧取回真实拍到的信息。这要求两件事同时成立：邻帧确实提供了不同的采样位置，且网络能把它们准确对齐到当前帧。

先看邻帧为什么能提供新信息。考虑一个简化的一维清晰信号$\bm{x}=(x_1,x_2,x_3,x_4)^\top$。第一帧只采到$x_1,x_3$，第二帧因半个低分辨率像素的位移采到$x_2,x_4$。在位置精确、没有模糊和噪声的教学条件下，两帧合起来给出全部四个分量，而单帧只给出两个。真实相机还进行光学模糊和像素积分，观测不是任意选出的几个高分辨率点；遮挡、光流偏差和颜色变化还会破坏互补条件。多帧超分可以利用这种观测差异，却不能宣称所有新增像素都被真实拍到。

在BasicVSR之前，视频超分常用滑动窗口：每个输出帧读取前后几帧，用复杂的对齐与融合模块处理。BasicVSR的作者把视频超分拆成传播、对齐、聚合、上采样四个功能，逐项检查哪种最简单的选择已经足够\scite{vision-basicvsr2021}。他们观察到两点：窗口方法看不到远处帧，而缩短可用范围会降低质量；单向循环方法中，序列前部的帧得到的信息少，质量随之下降。这两点分别指向“看得更远”和“前后都看”，所以BasicVSR的主要想法是用双向循环传播取得整段信息，再配以光流特征对齐、特征拼接聚合和像素重排上采样这几项简单选择。作者的结论是传播和对齐对质量与效率影响最大。

BasicVSR的输入是低分辨率序列$\bm{Y}_1,\dots,\bm{Y}_n$，每帧$h\times w\times3$。光流网络SPyNet（结构见\ref{sec:vision-video-warping}节）以预训练权重初始化，估计每对相邻帧之间两个方向的光流，每份形状为$h\times w\times2$。以前向分支为例，在时刻$t$，它把前一状态$\bm{h}^{\to}_{t-1}$按光流$\bm{u}_{t\to t-1}$反向取样到当前坐标，再与当前低分辨率帧一起送入由30个残差块组成的更新网络，得到$\bm{h}^{\to}_t$；状态有64个通道，形状为$h\times w\times64$。后向分支对称地读取$t+1$。用$\mathcal{W}$表示反向取样、$R_{\to}$和$R_{\leftarrow}$表示两路更新网络：
\begin{equation}
  \begin{aligned}
    \bm{h}^{\to}_t&=
      R_{\to}\!\left(\bm{Y}_t,\,
      \mathcal{W}(\bm{h}^{\to}_{t-1},\bm{u}_{t\to t-1})\right),\\
    \bm{h}^{\leftarrow}_t&=
      R_{\leftarrow}\!\left(\bm{Y}_t,\,
      \mathcal{W}(\bm{h}^{\leftarrow}_{t+1},\bm{u}_{t\to t+1})\right)\text{。}
  \end{aligned}
  \label{eq:vision-video-basicvsr}
\end{equation}
序列两端缺少的状态初始化为零。两路的参数各自在所有时刻共享，循环状态的一般性质可参考\BookChapterRef{chap:recurrent-networks}{卷五的循环网络章}。这里的状态不经门控，直接由残差块覆盖更新；能保留多少历史取决于训练学到的变换。原文只把更新网络称为“残差块的堆叠”，没有给出当前帧与取样状态在入口处怎样合并，也没有给出块内的层序与卷积核大小；读者复现时应以作者公开的实现为准，本书不补写这些细节。

聚合与上采样在每个时刻独立进行。把$\bm{h}^{\to}_t$与$\bm{h}^{\leftarrow}_t$沿通道拼接得到$h\times w\times128$的特征，再由若干卷积和\term{像素重排}（pixel shuffle）放大。像素重排是一种无参数的重新排列：把$h\times w\times(r^2c)$的特征中每个位置的$r^2$组通道，分别放到放大后$r\times r$小块的各个位置，得到$rh\times rw\times c$。例如$r=2$、$c=1$时，某位置的4个通道值$(v_1,v_2,v_3,v_4)$依次成为放大网格中该$2\times2$小块的左上、右上、左下、右下四个像素。它只改变排列，新增高分辨率值由前面卷积学到的通道内容决定。BasicVSR做4倍放大，原文说上采样模块由多个卷积和像素重排组成，输出$4h\times4w\times3$的画面，具体层数与通道没有在论文中列出。表\ref{tab:vision-video-basicvsr-layers}汇总论文给出的结构。

\begin{table*}[tbp]
  \centering
  \small
  \begin{tabularx}{\linewidth}{@{}>{\raggedright\arraybackslash}p{22mm}>{\raggedright\arraybackslash}X>{\raggedright\arraybackslash}X>{\raggedright\arraybackslash}p{17mm}>{\raggedright\arraybackslash}X>{\raggedright\arraybackslash}X@{}}
    \toprule
    阶段或层组 & 层类型 & 卷积核与步长 & 输出通道 & 重复次数 & 输出尺寸\\
    \midrule
    输入 & — & — & 3 & — & $h\times w$，每帧\\
    光流估计（SPyNet，每级） & 卷积，除末层外接ReLU & $7\times7$，步长1 & 32, 64, 32, 16, 2 & 每级5层，共5级 & $h\times w\times2$，每对相邻帧两个方向\\
    前向传播：对齐 & 双线性反向取样 & — & 64 & 每帧1次 & $h\times w\times64$\\
    前向传播：更新 & 残差块 & 论文未给出 & 64 & 30 & $h\times w\times64$\\
    后向传播：对齐与更新 & 同前向，参数独立 & 论文未给出 & 64 & 30 & $h\times w\times64$\\
    聚合 & 两路状态沿通道拼接 & — & 128 & 1 & $h\times w\times128$\\
    上采样 & 卷积与像素重排 & 论文未给出 & 3 & 论文未给出 & $4h\times4w\times3$\\
    \bottomrule
  \end{tabularx}
  \caption{BasicVSR的结构，按原论文正文与附录整理，输入为4倍超分辨率的低分辨率帧。SPyNet的层配置取自SPyNet原论文，BasicVSR未说明对它作了改动。原文给出的参数量为：光流网络140万，其余部分490万，合计630万；在$180\times320$的低分辨率输入上，表中报告的运行时间为63毫秒。}
  \label{tab:vision-video-basicvsr-layers}
\end{table*}

训练使用低分辨率序列与高分辨率目标，用式\eqref{eq:vision-video-charbonnier}的Charbonnier损失约束全部输出帧，BasicVSR取$\epsilon=10^{-8}$。低分辨率输入由两种规定的退化生成：双三次下采样，以及先做标准差1.6的高斯模糊再每隔四个像素取样；两种退化分别训练和评价。在REDS上每个样本取15帧并对15个输出都计算损失，在Vimeo-90K上把序列翻转拼接成14帧以加长传播；低分辨率块$64\times64$，批量8，用Adam和余弦退火训练30万步，主体学习率$2\times10^{-4}$。光流网络以预训练权重初始化，前5000步冻结，之后以$2.5\times10^{-5}$的较小学习率一起更新；它不使用光流真值，而是通过重建损失学习对超分有用的对应。推理时把整段视频一次输入。同一论文还提出IconVSR，在每隔五帧的关键帧上用额外特征提取器补充信息，并把后向分支特征送入前向分支；这些是对BasicVSR的扩展，不属于上述基本结构。

三帧车牌例子可以检查每个组件的作用。第一帧清楚的左侧笔画经前向状态传到第二帧，第三帧清楚的右侧笔画经后向状态补充第二帧。去掉对齐，同一位置可能来自不同笔画；去掉传播，只剩当前帧；去掉上采样，输出网格不会增加。拼接只让两路信息同时可用，是否被正确利用仍由训练和输入条件决定。图\ref{fig:vision-basicvsr-propagation}把两条时间依赖与一个前向更新分别画出，可以沿取样、合并和更新检查历史怎样进入当前状态。

\begin{figure}[htbp]
  \centering
  % 上半图省略每列共有的输入/重建边，以突出两条时间依赖。
\begin{tikzpicture}[x=1cm,y=1cm,font=\figurefont\small,
  state/.style={draw=BookTeal,fill=BookTeal!10,minimum width=17mm,
    minimum height=7mm,line width=0.8pt},
  block/.style={draw=BookInk,fill=BookPaper,minimum width=18mm,
    minimum height=8mm,align=center,line width=0.8pt},
  arr/.style={-{Stealth[length=2mm]},BookInk,line width=0.8pt}]
  \node[anchor=west] at (0,4.8) {两条独立的时间传播};
  \foreach \x/\i in {1.3/1,5/2,8.7/3}{
    \node at (\x,4.2) {时刻$\i$};
    \node[state] (f\i) at (\x,3.5) {$\bm h^{\to}_{\i}$};
    \node[state] (b\i) at (\x,2.4) {$\bm h^{\leftarrow}_{\i}$};
  }
  \draw[arr] (f1) -- (f2);
  \draw[arr] (f2) -- (f3);
  \draw[arr] (b3) -- (b2);
  \draw[arr] (b2) -- (b1);
  \node[anchor=west] at (0,1.45) {一个前向更新};
  \node[block] (old) at (1.1,0.6) {历史状态\\$\bm h^\to_{t-1}$};
  \node[block] (warp) at (4.1,0.6) {取样对齐};
  \node[block] (concat) at (7.0,0.6) {与当前帧合并};
  \node[state] (new) at (9.4,0.6) {$\bm h^\to_t$};
  \draw[arr] (old) -- (warp);
  \draw[arr] (warp) -- (concat);
  \draw[arr] (concat) -- node[above] {$R_\to$} (new);
  \node (flow) at (4.1,-0.6) {$\bm u_{t\to t-1}$};
  \node (input) at (7.0,-0.6) {$\bm Y_t$};
  \draw[arr] (flow) -- (warp);
  \draw[arr] (input) -- (concat);
\end{tikzpicture}

  \caption{BasicVSR的时间依赖与单步更新。上半图只画两路状态传播，每列都使用该时刻输入帧，同一列两路状态共同用于重建。下半图展开一个前向更新：历史状态按光流取样，再与当前帧一起送入残差块更新。后向分支必须等未来帧可用。}
  \label{fig:vision-basicvsr-propagation}
\end{figure}
\FloatBarrier

双向方法依赖未来帧。若改成单向在线运行，结构、训练和边界条件都发生变化，应重新训练和验证，不能只删除后向结果就沿用原成绩。大运动、遮挡和切镜可能使历史特征被拖入错误区域；车牌即使看上去变清晰，也须回查字符是否有观测支持。

BasicVSR与RVRT可以在同一任务（4倍视频超分、相同退化与评价协议）下比较，表\ref{tab:vision-video-rvrt-basicvsr}列出二者的主要设计。RVRT保留了BasicVSR的光流引导与循环传播，改变了三处：把逐帧更新换成$N$帧片段的联合更新，把单位置特征取样换成光流引导下的多位置注意力取样，并用多层交替方向代替两条独立分支。换来的是片段内帧间的直接信息交换和对光流误差更强的容忍；付出的代价是每步计算更重，参数量从630万增至1080万，训练所用帧数和步数也更多。RVRT原文在$320\times180$输入上的对比中，还列出了BasicVSR的后续版本BasicVSR++（730万参数、77毫秒），RVRT为183毫秒；作者认为运行时间的差距主要来自卷积层在现有深度学习框架中优化得更充分。两者都要等待整段输入，都不是在线方法。

\begin{table}[htbp]
  \centering
  \begin{tabularx}{\linewidth}{@{}lXX@{}}
    \toprule
    设计维度 & BasicVSR & RVRT\\
    \midrule
    传播单位 & 单帧状态，两条独立方向分支 & $N$帧片段，$L$层交替方向\\
    对齐方式 & 沿光流单位置反向取样特征 & 光流引导、每帧$M$个偏移位置的注意力取样\\
    帧内细化 & 卷积残差块 & 时空窗口内的Transformer块\\
    聚合方式 & 两路状态沿通道拼接 & 对齐特征与历史各层特征拼接后由注意力细化\\
    参数量 & 630万 & 1080万\\
    在线性 & 需整段，非因果 & 需整段，非因果\\
    \bottomrule
  \end{tabularx}
  \caption{同一视频超分任务下BasicVSR与RVRT的设计差别。两者都用预训练的SPyNet初始化并以Charbonnier损失端到端训练；参数量取自各自原论文，表中只列结构差别，不表示在任意数据上的质量排序。}
  \label{tab:vision-video-rvrt-basicvsr}
\end{table}

\subsection{中间时刻画面的恢复}
\label{sec:vision-video-interpolation}
\label{sec:vision-video-film}

\term{视频插帧}（video frame interpolation）根据两个时刻的观测，估计两者之间某一时刻的画面。输入是两帧$\bm{I}_0,\bm{I}_1$，形状$H\times W\times3$；输出是时刻$\tau\in(0,1)$的画面$\widehat{\bm{I}}_\tau$。它有两侧约束；只依据历史预测未来缺少后一侧约束，属于不同设置。缺帧补全也要注明缺失长度和实际可用的邻帧，缺一帧与缺半秒难度不同。

插帧的难点在于目标时刻没有观测，网络既要推算每个像素在$\tau$时刻的位置，又要在两侧都看不到的区域补出内容。经典做法沿用运动补偿的思路：估计两帧之间的光流，假设匀速运动把它按$\tau$缩放，从两侧反向取样后按时间距离加权平均。这一做法在运动边界和遮挡处会产生重影，因为缩放后的光流定义在错误的网格上，被遮挡的一侧也照样参与平均。基于学习的光流路线先估计运动，再从两侧反向取样并由网络合成；此前常见做法用一个预训练光流或深度网络估计运动，再用另一个网络合成，这需要稀缺的光流或深度真值，并且两辆车在两帧之间移动很多像素时，单一尺度的局部匹配可能找不到同一边缘\scite{vision-film2022}。

FILM（frame interpolation for large motion）的主要想法是：用跨尺度共享权重的特征提取器，让粗尺度上的大运动与细尺度上的小运动用同一套特征表示；在此基础上从粗到细直接预测从中间时刻指向两侧输入的位移，整个网络只用帧三元组训练。作者指出，标准的特征金字塔有两个困难：快速移动的小物体在粗层消失，粗层像素少、能提供的大运动监督也少。共享权重之所以有帮助，是因为大位移在粗尺度上变成小位移，而训练数据中小位移样本充足；同一组卷积学会处理小位移后，在粗尺度上就能处理原图中的大位移。

FILM的特征提取先把每张输入建成7层图像金字塔，第$l$层分辨率是原图的$2^{-(l-1)}$。每层图像都送入同一个编码器，编码器第$d$个深度的输出记为$\bm{f}^{l,d}$，原实现共4个深度；每个深度是两层$3\times3$卷积加ReLU，之后接一个大小和步长都为2的平均池化，相同深度在所有金字塔层共享卷积权重。于是第$l$层图像在深度2的特征与第$l+1$层图像在深度1的特征具有相同空间分辨率。FILM把相同分辨率、来自不同深度和不同层的特征沿通道拼接，记为$\bm{F}^{l}=(\bm{f}^{l-2,3},\bm{f}^{l-1,2},\bm{f}^{l,1})$；最细的两层可用的项较少，分别只拼接一份和两份。这样，每个尺度上的特征都混合了“浅层看细节”和“深层看大范围”的信息，含义在各尺度之间可比，原文称之为尺度无关的特征。

运动估计从最粗层开始。第$l$层要预测从中间时刻指向第1帧的位移$\bm{w}^{l}_{\tau\to1}$：先把上一粗层结果放大到本层网格，记为$(\bm{w}^{l+1}_{\tau\to1})_{\times2}$，用它把第1帧的特征$\bm{F}_1^{l}$反向取样，再由一个小卷积网络$G^l$读取第0帧特征和取样结果，预测残差：
\begin{equation}
  \bm{w}^{l}_{\tau\to1}
   =(\bm{w}^{l+1}_{\tau\to1})_{\times2}
   +G^{l}\Bigl(\bm{F}_0^{l},\,
     \mathcal{W}\bigl(\bm{F}_1^{l},(\bm{w}^{l+1}_{\tau\to1})_{\times2}\bigr)\Bigr)\text{。}
  \label{eq:vision-video-film-flow}
\end{equation}
这与式\eqref{eq:vision-video-spynet}的SPyNet更新形式相同，差别在于FILM在共享特征上而不是在图像上做残差估计，位移的起点也换成了中间时刻。原文把$(\cdot)_{\times2}$写作两倍双线性上采样；位移若以当前网格的像素为单位，网格放大一倍时数值也应随之加倍，读者复现时要确认实现采用的单位约定。指向第0帧的位移$\bm{w}^{l}_{\tau\to0}$对称计算。原文结构图中，$G^l$为三层$3\times3$卷积加ReLU，在较粗的第3至7层共享权重，最细的两层各用自己的参数。粗尺度上原图40像素的运动在第4层只有5像素，较容易进入局部比较范围，再逐层加上细的修正。这些位移由最终中间帧的监督间接决定，原文称为面向任务的位移，它不必等于真实光流，也不能当作独立验证的光流使用。

以教学数值说明两路取样方向。设$\tau=1/2$，一辆车在两端横坐标20和28之间匀速平移，中间时刻位置为24。输出位置24应分别向第0帧的20和第1帧的28取样，位移为$-4$和$+4$。真实模型还要处理加速、遮挡和新显露的背景，不是对两个输入坐标机械取平均；两侧都没见过的区域只能由合成网络依先验估计。

合成阶段先用各层的位移把两侧的特征与该层图像一起反向取样到中间时刻，再在每个金字塔层把两侧取样结果和两份位移拼接，送入类似U-Net的解码器。解码器从粗到细逐层上采样，与本层的拼接输入合并，经两层卷积加ReLU，最后一层卷积输出$H\times W\times3$的中间帧。表\ref{tab:vision-video-film-layers}汇总论文给出的结构。原文报告，在V100 GPU上插一帧720p画面约0.39秒，峰值显存约4.5 GB；论文没有给出参数量。

\begin{table*}[tbp]
  \centering
  \small
  \begin{tabularx}{\linewidth}{@{}>{\raggedright\arraybackslash}p{22mm}>{\raggedright\arraybackslash}X>{\raggedright\arraybackslash}X>{\raggedright\arraybackslash}p{17mm}>{\raggedright\arraybackslash}X>{\raggedright\arraybackslash}X@{}}
    \toprule
    阶段或层组 & 层类型 & 卷积核与步长 & 输出通道 & 重复次数 & 输出尺寸\\
    \midrule
    输入 & 两帧 & — & 3 & 2 & $H\times W$\\
    图像金字塔 & 逐级缩小 & — & 3 & 7级 & 第$l$级$H/2^{l-1}\times W/2^{l-1}$\\
    共享编码器（每个深度） & 卷积加ReLU，再平均池化 & $3\times3$，步长1；池化$2\times2$，步长2 & 第1深度64，其余论文未给出 & 每深度2层卷积，共4个深度 & 每个深度减半\\
    尺度无关特征 & 同分辨率特征拼接 & — & 三个深度通道之和 & 每级1次 & 第$l$级分辨率\\
    残差位移估计$G^l$ & 卷积加ReLU & $3\times3$，步长1 & 2（每个方向） & 3层；第3至7级共享 & 第$l$级分辨率\\
    融合解码器（每级） & 上采样、拼接、卷积加ReLU & 论文未给出 & 论文未给出 & 每级2层 & 逐级放大\\
    输出层 & 卷积 & 论文未给出 & 3 & 1 & $H\times W\times3$\\
    \bottomrule
  \end{tabularx}
  \caption{FILM的结构，按原论文正文和结构图整理。论文正文说明所用模型的编码器从64个滤波器开始，结构图以32起始作示意，未逐项列出更深层的通道数；结构图标注的$G^l$为三层带ReLU的卷积，末层输出两通道位移。论文未给出参数量。}
  \label{tab:vision-video-film-layers}
\end{table*}

训练样本是两端与中间真实帧组成的三元组，只在$\tau=1/2$监督。损失作用在最终输出上：
\begin{equation}
  L_{\mathrm{FILM}}
  =w_{1}\lVert\widehat{\bm{I}}_{1/2}-\bm{I}_{1/2}\rVert_1
   +w_{\mathrm{VGG}}L_{\mathrm{VGG}}
   +w_{\mathrm{Gram}}L_{\mathrm{Gram}}\text{，}
  \label{eq:vision-video-film-loss}
\end{equation}
第一项是像素绝对误差，约束颜色和位置；$L_{\mathrm{VGG}}$比较预训练VGG-19网络各层特征的绝对误差，约束局部结构；$L_{\mathrm{Gram}}$比较这些特征的Gram矩阵（通道两两相关），约束纹理统计，使大运动区域不至于被平均得过于模糊。后两项的定义和作用见\ref{sec:vision-generation-perceptual-style}节。原实现在前150万步取$(w_1,w_{\mathrm{VGG}},w_{\mathrm{Gram}})=(1,1,0)$，后150万步转为$(1,0.25,40)$，使三项贡献大致相当；原文还训练了只用第一项和只加VGG项的版本：各版本在像素指标上差距很小，但原文展示只用第一项时画面偏模糊，加入Gram项后明显变锐，报告结果时应注明使用哪一版本。训练使用Vimeo-90K三元组，$256\times256$随机裁剪，批量8，用Adam优化共300万步，学习率从$10^{-4}$开始每75万步按0.464的比例衰减，并以旋转、翻转和三元组时间反转增强。原文还从视频中挖掘不同运动幅度范围的三元组，报告用较大运动范围的数据训练有助于大运动测试集。

需要更密的时间网格时，FILM递归插入中点，得到$1/4,1/2,3/4$等位置。新一轮可能把先前生成的帧作为输入，误差会继承，生成帧不能当作新增的实拍约束。在线插帧必须等待后一端：30帧每秒的输入估计中点时，中点内容至少要在后一帧到达后才能生成，相对该中点的观测等待约为16.7毫秒，还要加上计算耗时和播放缓冲。平滑慢动作适合呈现，却不能证明插出的碰撞、接触或遮挡细节真实发生过。

FILM与BasicVSR都用反向取样和多尺度运动，但目的不同，可以从三个方面比较。第一，取样起点：BasicVSR从已有帧的网格出发读取邻帧，FILM从一个尚无观测的时刻出发读取两侧。第二，运动来源：BasicVSR用预训练的SPyNet初始化后微调，FILM的运动估计完全由帧监督学到，不依赖光流或深度真值。第三，输出依据：BasicVSR的每个输出位置至少有本帧的低分辨率观测，FILM的输出完全是由两侧推算的估计。FILM因此省去了光流数据，但它的位移和输出都缺少独立的观测检查，只能靠两侧一致性和下游用途核验。

\section{有限比特下的视频编码与解码}
\label{sec:vision-video-codec}

前两节利用帧间对应改善画面，本节利用同一对应减少传输。问题是：输入原视频$\bm{Y}_1,\dots,\bm{Y}_n$，输出一条码流，使解码端只凭码流就能重建$\widehat{\bm{Y}}_t$，并在给定比特预算下让失真尽量小。相邻帧有大量相同内容，逐帧独立编码会反复传输它们；帧间编码改为用已解码的帧预测当前帧，只编码预测解释不了的部分。这里待编码的输入仍记为$\bm{Y}_t$，压缩目标是保持这个实际输入，不默认它已是前两节追求的无退化真值；压缩也可以直接处理原视频，不要求先复原。

难点在于编码器与解码器看到的信息不同。编码器能看原视频，解码器只能读码流和自己已经重建的帧；任何预测都必须只用两端共有的信息，否则两端的重建会逐帧分叉。此外，运动和残差都要花比特，减少一方可能增加另一方，必须联合权衡。

视频编码标准对这些难点的回答是\term{混合视频编码}（hybrid video coding）：用帧间或帧内预测去掉时间和空间上可预测的部分，再对预测残差做二维变换编码。从H.261起的各代视频编码标准都采用这一结构，H.264/AVC和HEVC也不例外\scite{vision-sullivan2012hevc}。图\ref{fig:vision-video-hybrid-codec}按编码一帧的顺序画出主要模块。编码器把图像划分成块；序列首帧和每个随机访问点只用帧内预测，其余帧的多数块用帧间预测，编码器为每块选定参考帧和运动矢量，作为辅助信息传给解码器。预测残差经变换、缩放、量化和熵编码写入码流。编码器内部复制一遍解码器的处理：把量化后的系数反量化、反变换，加回预测，经环路滤波后存入已解码帧缓存，供后续帧预测使用；这样两端才能产生相同的预测。标准只规定码流语法和解码过程，编码器怎样搜索运动、怎样选择模式都由实现自行决定\scite{vision-wiegand2003h264}。

\begin{figure*}[tbp]
  \centering
  % H.264/HEVC式混合编码器的主要模块，按两篇标准概述的编码流程重绘；宽约150mm，跨栏使用。
% 灰底模块组成编码器内部的解码环路，与独立解码器执行相同的计算；白底模块只在编码端出现。
\begin{tikzpicture}[x=1cm,y=1cm,font=\figurefont\small,
  enc/.style={draw=FigureBlue,fill=white,rounded corners=1.8pt,line width=0.8pt,
    minimum width=21mm,minimum height=8mm,align=center,inner sep=1.2mm},
  dec/.style={enc,draw=FigureStroke,fill=FigurePaper},
  sum/.style={circle,draw=FigureStroke,fill=white,line width=0.8pt,
    minimum size=5.5mm,inner sep=0pt},
  arr/.style={-{Stealth[length=1.8mm]},draw=FigureStroke,line width=0.8pt,rounded corners=1mm},
  side/.style={arr,dashed},
  note/.style={text=FigureMuted,font=\figurefont\footnotesize}]
  \node[enc] (in) at (0,3) {当前帧\\$\bm Y_t$};
  \node[sum] (sub) at (2.4,3) {$-$};
  \node[enc] (tr) at (4.6,3) {整数变换};
  \node[enc] (q) at (7.2,3) {量化\\（量化参数）};
  \node[enc] (ent) at (10.6,3) {熵编码\\CAVLC/CABAC};
  \node[enc,minimum width=14mm] (bs) at (14.0,3) {码流};
  \node[dec] (iq) at (7.2,1.5) {反量化\\反变换};
  \node[sum] (add) at (7.2,0.2) {$+$};
  \node[dec] (lf) at (9.8,0.2) {环路滤波\\去块、SAO};
  \node[dec] (buf) at (12.6,0.2) {已解码帧\\缓存};
  \node[enc] (me) at (12.6,1.7) {运动估计};
  \node[dec] (pred) at (2.4,-1.3) {帧内预测或\\运动补偿};
  \draw[arr] (in) -- (sub);
  \draw[arr] (sub) -- node[above,note] {残差} (tr);
  \draw[arr] (tr) -- (q);
  \draw[arr] (q) -- node[above,note] {系数} (ent);
  \draw[arr] (ent) -- (bs);
  \draw[arr] (q) -- (iq);
  \draw[arr] (iq) -- (add);
  \draw[arr] (add) -- (lf);
  \draw[arr] (lf) -- (buf);
  \draw[arr] (buf) -- (me);
  \draw[arr] (in.north) -- (0,4.1) -| (me.north);
  \draw[side] (me.west) -| node[pos=0.75,left,note] {运动矢量、模式} (ent.south);
  \draw[arr] (buf.south) |- node[pos=0.75,above,note] {参考帧} (pred.east);
  \draw[arr] (pred.north) -- node[left,note] {预测$\overline{\bm Y}_t$} (sub.south);
  \draw[arr] (2.4,0.2) -- (add.west);
  \draw[side] (me.east) -- (14.1,1.7) |- (2.4,-2.3) -- (pred.south);
  \node[note,anchor=west] at (3.0,-2.6) {灰底：编码器内部的解码环路，与解码器逐位一致};
\end{tikzpicture}

  \caption{H.264/HEVC式混合编码器的主要模块，根据两篇标准概述的编码流程重绘。实线为数据流，虚线为运动矢量和模式等辅助信息。灰底模块构成编码器内部的解码环路，独立解码器执行同样的计算；运动估计和模式选择只在编码端进行，不属于标准规定的范围。}
  \label{fig:vision-video-hybrid-codec}
\end{figure*}

这套框架的每个模块都按手工设计分别调优：块匹配找运动，固定的整数变换去相关，量化参数控制精度，熵编码器按上下文切换概率，编码器再用率失真代价为每个块选择模式。图像压缩的两条路线已在\ref{sec:vision-restoration-compression}节讲过：JPEG式的经典编码（\ref{sec:vision-restoration-classic-codec}节）用块变换、量化表和熵编码处理单幅图像，学习式编码（\ref{sec:vision-restoration-learned-codec}节）证明了变换、量化和概率模型可以联合训练，但二者都不处理帧间依赖。本节只讲视频新增的部分：帧间预测与参考结构、整数变换与量化参数、CABAC、环路滤波和模式选择。DVC沿用传统的“运动补偿—残差编码”框架，把每个模块换成网络，用一个率失真目标联合训练全部模块\scite{vision-dvc2019}。表\ref{tab:vision-video-codec-modules}按模块对照两者。编码流程中的三个步骤分别面对上面的难点：确定参考依赖解决两端信息不对称，运动与残差的编码解决比特分配，解码重建检验依赖能否在接收端独立还原。每一步先讲H.264/HEVC的做法，再讲DVC改了什么。

\begin{table}[htbp]
  \centering
  \small
  \begin{tabularx}{\linewidth}{@{}lXX@{}}
    \toprule
    模块 & H.264/HEVC & DVC\\
    \midrule
    运动估计 & 编码器自选的块搜索；可变块大小，四分之一像素精度 & 预训练后联合微调的SPyNet，逐像素光流\\
    运动信息编码 & 运动矢量由相邻块预测，差值熵编码 & 卷积自编码器压缩光流，学习的概率模型估计码率\\
    运动补偿 & 按块从参考帧复制，分数位置插值 & 双线性反向取样后由补偿网络修正\\
    残差变换 & 近似离散余弦变换的整数变换 & 尺度超先验方法的非线性变换\\
    量化 & 由量化参数控制的均匀量化 & 取整，训练时以加均匀噪声代替\\
    熵编码 & CAVLC或CABAC & 按学习的符号概率熵编码\\
    环路滤波 & 去块滤波；HEVC另有样点自适应偏移 & 不设，原文认为逐像素补偿不产生块效应\\
    模式与比特分配 & 编码时逐块最小化拉格朗日代价 & 训练时最小化率失真损失，每个$\lambda$一个模型\\
    \bottomrule
  \end{tabularx}
  \caption{混合编码框架与DVC的模块对照。两者的数据流相同，差别在于各模块由手工设计还是由网络实现，以及率失真权衡发生在编码时还是训练时。H.264/HEVC一列根据两篇标准概述整理，只列与本章相关的工具。}
  \label{tab:vision-video-codec-modules}
\end{table}

\subsection{参考依赖的确定}
\label{sec:vision-video-codec-reference}
\label{sec:vision-video-dvc}

编码第一步要决定每帧依赖什么。\term{帧内编码}（intra-frame coding）只用当前帧自身的信息，可以独立解码，作为序列起点或错误恢复点；\term{帧间编码}（inter-frame coding）依赖其他已重建帧，比特更少，却必须先解码它所依赖的帧。从一个帧内编码帧开始、到下一个帧内编码帧之前的一组帧称为图像组（group of pictures, GOP，见\ref{sec:vision-videofound-formats}节），它规定了参考链的长度。这一步的难点不是选择哪种编码，而是保证参考对两端相同，并控制依赖带来的等待。

H.264把一帧划分为若干条带，按可用的预测方式分为三类\scite{vision-wiegand2003h264}：I条带中的块全部用帧内预测；P条带中的块还可以用帧间预测，每个预测块至多有一个运动补偿预测信号；B条带中的块可以同时用两个运动补偿预测信号，例如一个来自显示顺序上更早的帧、一个来自更晚的帧。H.264还允许从多个已解码帧中选择参考，并把参考顺序与显示顺序解耦，B帧本身也可以作为其他帧的参考。下文按习惯把主要由这三类条带组成的帧称为I帧、P帧和B帧。

参考更晚的帧，意味着解码顺序与显示顺序不同。图\ref{fig:vision-video-gop-order}是一个七帧的教学图像组：P$_3$参考I$_0$，B$_1$与B$_2$同时参考I$_0$和P$_3$。要解码B$_1$，P$_3$必须先到达，所以码流中的顺序是I$_0$、P$_3$、B$_1$、B$_2$、P$_6$、B$_4$、B$_5$。编码端在拍到第3帧之前无法编码B$_1$，按每秒30帧计算，B$_1$至少多等两帧，约66.7毫秒；B帧能借两侧参考减少残差，代价正是\ref{sec:vision-video-restoration-availability}节所说的等待。只用I帧和P帧的结构没有这种重排等待，后文的DVC就属于这一类。

\begin{figure}[htbp]
  \centering
  % 七帧教学图像组：显示顺序与解码顺序；宽约100mm。实际编码器可选择不同的参考结构。
\begin{tikzpicture}[x=1cm,y=1cm,font=\figurefont\small,
  frame/.style={draw=#1,fill=#1Fill,rounded corners=1.8pt,line width=0.8pt,
    minimum width=9mm,minimum height=8mm,inner sep=0.8mm},
  ref/.style={-{Stealth[length=1.6mm]},draw=FigureStroke,line width=0.7pt},
  bref/.style={ref,dashed,draw=FigureMuted},
  note/.style={text=FigureMuted,font=\figurefont\footnotesize}]
  \node[anchor=east] at (-0.6,2.2) {显示顺序};
  \foreach \i/\t/\c in {0/I/FigureBlue,1/B/FigureRed,2/B/FigureRed,3/P/FigureYellow,
                        4/B/FigureRed,5/B/FigureRed,6/P/FigureYellow}{
    \node[frame=\c] (d\i) at ({1.2*\i},2.2) {$\mathrm{\t}_{\i}$};
  }
  \draw[ref] (d0.north) to[bend left=35] node[above,note] {参考} (d3.north);
  \draw[ref] (d3.north) to[bend left=35] (d6.north);
  \draw[bref] (d0.south) to[bend right=40] (d1.south);
  \draw[bref] (d3.south) to[bend left=40] (d1.south);
  \draw[bref] (d3.south) to[bend right=40] (d4.south);
  \draw[bref] (d6.south) to[bend left=40] (d4.south);
  \node[anchor=east] at (-0.6,0) {解码顺序};
  \foreach \k/\i/\t/\c in {0/0/I/FigureBlue,1/3/P/FigureYellow,2/1/B/FigureRed,3/2/B/FigureRed,
                           4/6/P/FigureYellow,5/4/B/FigureRed,6/5/B/FigureRed}{
    \node[frame=\c] (e\k) at ({1.2*\k},0) {$\mathrm{\t}_{\i}$};
  }
  \foreach \k in {0,...,5}{
    \pgfmathtruncatemacro{\n}{\k+1}
    \draw[ref,draw=FigureMuted] (e\k.east) -- (e\n.west);
  }
\end{tikzpicture}

  \caption{七帧教学图像组的参考关系与解码顺序。下标为显示时刻；上方实线箭头从参考帧指向P帧，下方虚线箭头为B$_1$和B$_4$的前后两个参考，B$_2$、B$_5$同理而未画出。解码顺序把每个B帧排在它的两个参考之后。这是教学构造，实际编码器可以选择不同的参考结构。}
  \label{fig:vision-video-gop-order}
\end{figure}

参考对两端相同，不只要求两端用同一帧，还要求两端对这一帧的重建逐位一致。早期标准只在一个误差容限内规定反变换，不同解码器的输出略有差异，这种差异会沿参考链累积，称为漂移；H.264把反变换定义为精确的整数运算，成为第一个使所有解码器输出完全相同的标准\scite{vision-wiegand2003h264}。这个要求同样适用于学习式编码，后文会看到它给神经网络带来的新问题。

DVC在帧间编码中只用上一重建帧作参考，即只有前向预测的低延迟结构。它的动机有两点：传统编码器的块运动估计和离散余弦变换是手工设计，整个系统没有端到端优化；此前的学习式视频压缩或只替换单个模块，或仍依赖传统块运动估计，或只优化失真而不计码率。DVC的主要想法是保留“估计运动、压缩运动、\term{运动补偿}（motion compensation）、压缩残差”这一被长期验证的结构，用网络实现每个模块，再用率失真损失共同训练。运动补偿是指用运动信息把参考帧变换成当前帧的预测；网络化之后，运动可以是逐像素的稠密光流而非块矢量，预测不会出现块边界。

DVC的数据流如下。输入当前帧$\bm{Y}_t$（$H\times W\times3$）和上一重建帧$\widehat{\bm{Y}}_{t-1}$。光流网络估计光流$\bm{v}_t$，形状$H\times W\times2$；运动编码器把它压缩为潜在表示$\bm{m}_t$，量化为$\widehat{\bm{m}}_t$，运动解码器再恢复出$\widehat{\bm{v}}_t$。运动补偿用$\widehat{\bm{v}}_t$对$\widehat{\bm{Y}}_{t-1}$反向取样，再由补偿网络修正，输出预测帧$\overline{\bm{Y}}_t$。残差$\bm{E}_t$经残差编码器变为潜在表示$\bm{z}_t$，量化为$\widehat{\bm{z}}_t$后解码为$\widehat{\bm{E}}_t$：
\begin{equation}
  \bm{E}_t=\bm{Y}_t-\overline{\bm{Y}}_t,\qquad
  \widehat{\bm{Y}}_t=\overline{\bm{Y}}_t+\widehat{\bm{E}}_t\text{。}
  \label{eq:vision-video-dvc-reconstruction}
\end{equation}
$\bm{E}_t$是编码端可算的残差，$\widehat{\bm{E}}_t$是它经编码、量化与解码后的重建。当前重建帧写入参考缓存，成为下一帧的共同条件。各网络的层结构在下一小节给出。

关键的约束是：运动补偿必须使用解码后的$\widehat{\bm{v}}_t$，而不是光流网络直接输出的$\bm{v}_t$。原始光流未经编码量化，解码端拿不到；若编码端用它产生预测，就算出了解码端无法复现的残差，也掩盖了运动压缩误差。图\ref{fig:vision-dvc-reference}按两端职责画出这条依赖，它与图\ref{fig:vision-video-hybrid-codec}中编码器内部的解码环路是同一个要求。

\begin{figure*}[tbp]
  \centering
  % 原图仅编码端可见；下路也由编码端本地执行，保证共同参考。
\begin{tikzpicture}[x=1cm,y=1cm,font=\figurefont\small,
  block/.style={draw=BookInk,fill=BookPaper,minimum width=24mm,
    minimum height=9mm,align=center,line width=0.8pt},
  arr/.style={-{Stealth[length=2mm]},BookInk,line width=0.8pt,rounded corners=1mm}]
  \node[anchor=west] at (-0.2,4.7) {编码端：从原帧与本地预测形成两类符号};
  \node[block] (srcm) at (1.4,3.6) {$\bm Y_t,\widehat{\bm Y}_{t-1}$};
  \node[block] (encm) at (5.2,3.6) {运动估计\\编码、量化};
  \node[block] (symm) at (9.0,3.6) {$\widehat{\bm m}_t$};
  \node[block] (srcr) at (1.4,2.2) {$\bm E_t=\bm Y_t-\overline{\bm Y}_t$};
  \node[block] (encr) at (5.2,2.2) {残差编码\\量化};
  \node[block] (symr) at (9.0,2.2) {$\widehat{\bm z}_t$};
  \draw[arr] (srcm) -- (encm);
  \draw[arr] (encm) -- (symm);
  \draw[arr] (srcr) -- (encr);
  \draw[arr] (encr) -- (symr);
  \draw[BookTeal,fill=BookTeal!8,line width=0.8pt] (12.7,1.65) rectangle (15.6,4.15);
  \node at (14.15,3.6) {运动码流};
  \node at (14.15,2.2) {残差码流};
  \draw[arr] (symm.east) -- node[above] {熵编码} (12.7,3.6);
  \draw[arr] (symr.east) -- node[above] {熵编码} (12.7,2.2);
  \draw[BookRule,line width=0.6pt] (-0.2,1.2) -- (15.6,1.2);
  \node[anchor=west] at (-0.2,0.6) {解码端：先从码流恢复符号};
  \node[block] (dm) at (1.4,-0.2) {$\widehat{\bm m}_t$};
  \node[block] (decm) at (5.2,-0.2) {运动解码};
  \node[block] (comp) at (9.0,-0.2) {运动补偿};
  \node[circle,draw=BookTeal,fill=BookTeal!8,line width=0.8pt,
    minimum size=6mm,inner sep=0pt] (sum) at (12.2,-0.2) {$+$};
  \node[block] (out) at (14.4,-0.2) {$\widehat{\bm Y}_t$};
  \draw[arr] (dm) -- (decm);
  \draw[arr] (decm) -- node[above] {$\widehat{\bm u}_t$} (comp);
  \draw[arr] (comp) -- node[above] {$\overline{\bm Y}_t$} (sum);
  \draw[arr] (sum) -- (out);
  \node (ref) at (9.0,0.7) {$\widehat{\bm Y}_{t-1}$};
  \draw[arr] (ref) -- (comp);
  \node[block] (dr) at (1.4,-1.7) {$\widehat{\bm z}_t$};
  \node[block] (decr) at (5.2,-1.7) {残差解码};
  \draw[arr] (dr) -- (decr);
  \draw[arr] (decr.east) -- node[above,pos=0.4] {$\widehat{\bm E}_t$}
    (12.2,-1.7) -- (sum.south);
  \node[align=center] (buffer) at (14.4,-1.7) {保存为\\下一帧参考};
  \draw[arr] (out.south) -- (buffer.north);
\end{tikzpicture}

  \caption{DVC式单参考帧间编码的主要依赖，根据原论文的计算职责重绘。上路在编码端形成量化运动与残差符号并写入码流，下路说明独立解码器如何从码流重建。两端运动补偿均使用解码后的运动和上一重建帧，编码器还须执行与下路相同的本地解码以更新参考；原始当前帧只在编码端可见。}
  \label{fig:vision-dvc-reference}
\end{figure*}
\FloatBarrier

用四个像素表示横向移动的方块。上一重建为$(0,1,1,0)$，向右移动一个像素、左边以零填充的预测为$(0,0,1,1)$。设当前源帧为$(0.2,0,1,1)$，左侧新进入的背景并非零，运动预测无法解释它，残差为$(0.2,0,0,0)$。若残差解码得到$(0.18,0,0,0)$，当前重建就是$(0.18,0,1,1)$，四像素均方误差为$0.02^2/4=10^{-4}$。这是省略了补偿网络修正的教学简化。

下一帧要使用含0.18的重建参考，而不能使用编码端掌握的原值0.2。否则，两端即使拿到相同的运动和残差，也会从不同起点产生不同图像，误差沿参考链继续传播。这就是编码器内部必须运行一遍本地解码的原因。首帧需要另行用帧内方式编码，DVC的帧间网络没有消除这个起点成本。

\subsection{运动与残差的编码}
\label{sec:vision-video-rate-allocation}

确定了参考，第二步要把运动与残差两类信息变成比特，并决定两者各花多少。输入是当前帧、参考帧和由此得到的运动与残差，输出是两段码流。难点在于运动越粗，运动比特越少，但预测越差，残差比特越多；只优化其中一方会把成本推给另一方。新显露区域等帧间预测失效的地方，还要决定是否改用帧内预测。

混合编码先确定在什么粒度上描述运动。H.264以\term{宏块}（macroblock）为基本单位，一个宏块覆盖$16\times16$的亮度样本和两个$8\times8$的色度块（4:2:0采样，见\ref{sec:vision-videofound-formats}节）\scite{vision-wiegand2003h264}。P宏块可以划分为$16\times16$、$16\times8$、$8\times16$或$8\times8$的预测块，$8\times8$块还可以再分成$8\times4$、$4\times8$或$4\times4$，于是一个宏块最多携带16个运动矢量。HEVC把宏块换成大小由编码器选择的\term{编码树单元}（coding tree unit, CTU），亮度块边长取16、32或64，再按四叉树递归划分成编码单元；预测块从$64\times64$到$4\times4$，变换块从$4\times4$到$32\times32$\scite{vision-sullivan2012hevc}。大块适合平坦背景，小块适合运动边界，划分本身也要花比特，这正是比特分配的第一个层面。

运动矢量由编码器搜索得到，常用\ref{sec:vision-videofound-block-matching}节的块匹配。H.264的运动矢量精度为四分之一像素：半像素位置用六抽头滤波器在水平和垂直方向插值，四分之一位置再由整数和半像素位置的样本平均得到；HEVC改用七抽头或八抽头滤波器。运动矢量本身先由相邻块的矢量预测，只传差值，HEVC还提供合并模式，直接继承相邻块的运动。新显露或运动复杂的区域可以改用帧内预测，用同一帧中已解码的相邻样本沿某个方向外推：H.264的$4\times4$帧内预测有9种模式、$16\times16$有4种，HEVC有33种方向模式加平面和直流两种模式。

预测残差经二维变换去掉空间相关。这一步与JPEG的$8\times8$离散余弦变换和量化表（\ref{sec:vision-restoration-jpeg-transform}节）作用相同，变换本身的定义见\ref{sec:vision-transforms-dct}节；视频标准的不同在于变换作用于预测残差而非像素，并改用整数近似。H.264主要使用$4\times4$的可分离整数变换，变换矩阵为
\begin{equation}
  \bm{T}=\begin{pmatrix}
    1&1&1&1\\
    2&1&-1&-2\\
    1&-1&-1&1\\
    1&-2&2&-1
  \end{pmatrix},\qquad
  \bm{C}=\bm{T}\bm{R}\bm{T}^{\top}\text{，}
  \label{eq:vision-video-h264-transform}
\end{equation}
其中$\bm{R}$是$4\times4$的残差块，$\bm{C}$是变换系数。$\bm{T}$的各行与$4\times4$离散余弦变换的基向量形状相近，但元素都是小整数，正反变换只需加法和移位；各行范数不同（2或$\sqrt{10}$），归一化并入量化时的缩放。一维教学例：一行线性增加的残差$(10,12,14,16)^\top$，乘以$\bm{T}$得到$(52,-14,0,-2)^\top$，能量几乎全集中在前两个系数，后两个取整后可为零。HEVC把整数变换扩展到$32\times32$，并对帧内亮度$4\times4$残差另行规定一种由离散正弦变换导出的整数变换。

变换系数按\term{量化参数}（quantization parameter, QP）均匀量化。H.264的量化参数有52个取值，每增加1，量化步长约增大12\%，每增加6，步长恰好加倍；按概述的经验，步长增大约12\%时码率约降低12\%\scite{vision-wiegand2003h264}。量化后的系数按之字形顺序扫描（与JPEG相同，见\ref{sec:vision-restoration-symbolization}节），再熵编码。H.264提供两种熵编码：\term{上下文自适应变长编码}（context-adaptive variable-length coding, CAVLC）按已传输的语法元素在多张变长码表之间切换，与JPEG用固定Huffman码表（\ref{sec:vision-restoration-huffman}节）相比多了上下文自适应；\term{上下文自适应二进制算术编码}（context-adaptive binary arithmetic coding, CABAC）先把每个语法元素二值化，再用由已编码元素构成的上下文选择概率模型，以算术编码输出比特\scite{vision-cabac2003}。算术编码的区间划分原理见\ref{sec:vision-restoration-entropy-coding}节；它能为一个符号分配非整数比特，概率还能随统计变化自适应。概述报告，CABAC相对CAVLC通常节省5\%到15\%的码率；HEVC只保留了改进后的CABAC。

这些模块都提供了多种选择：块怎样划分，用帧间还是帧内预测，参考哪一帧，运动矢量取哪个。编码器用\term{率失真优化}（rate-distortion optimization）在它们之间选择：对每个候选模式计算失真$D$和码率$R$，选使拉格朗日代价
\begin{equation}
  J=D+\lambda R
  \label{eq:vision-video-lagrangian}
\end{equation}
最小的那个\scite{vision-rdo-video1998}。$\lambda>0$是拉格朗日乘子，表示多花一比特愿意换回多少失真下降。这一做法的依据是约束优化中的乘子法：在码率预算下最小化失真，可以化为对某个$\lambda$求无约束的$D+\lambda R$的最小值，乘子的一般理论可参考\BookChapterRef{chap:optimization}{卷一的最优化章}。条件是所求的工作点位于全部可选（码率，失真）点的下凸包上；凸包以内的点无论取什么$\lambda$都选不到。由于各块的$D$和$R$可以相加，整帧的代价分解为逐块代价之和，于是可以逐块独立选择。率失真函数本身的信息论刻画卷一和卷二都没有展开，本章只使用这种可操作的形式。算法\ref{alg:vision-video-mode-decision}写出逐块选择的步骤。

\begin{algorithm}[基于拉格朗日代价的块模式选择]
  \label{alg:vision-video-mode-decision}
  \begin{algorithmic}[1]
    \Require 当前块$\bm{B}$；参考帧与已解码的相邻样本；候选模式集合$\mathcal{M}$；乘子$\lambda>0$
    \Ensure 所选模式$m^\star$
    \For{每个候选模式$m\in\mathcal{M}$}
      \State 按模式$m$形成预测$\overline{\bm{B}}_m$（帧内外推或按运动矢量补偿）
      \State 对残差$\bm{B}-\overline{\bm{B}}_m$做变换、量化和反变换，得到重建$\widehat{\bm{B}}_m$
      \State $D_m\gets\lVert\bm{B}-\widehat{\bm{B}}_m\rVert_2^2$；$R_m\gets$模式、运动矢量与残差系数的比特数
      \State $J_m\gets D_m+\lambda R_m$
    \EndFor
    \State $m^\star\gets\argmin_{m\in\mathcal{M}}J_m$
  \end{algorithmic}
\end{algorithm}

用一个$16\times16$块复算，数值为教学构造。四个候选的（失真，比特）分别为：跳过（直接复制参考，400，1）、单个$16\times16$运动矢量（120，30）、四个$8\times8$运动矢量（60，90）、帧内预测（50，150）。取$\lambda=4$，代价依次为404、240、420、750，选单个运动矢量；取$\lambda=0.5$，代价依次为400.5、135、105、125，改选四个$8\times8$运动矢量。$\lambda$变小意味着更看重失真，编码器愿意为更细的运动多花比特。这个$\lambda$与量化参数需要配套。在最优工作点附近，$\lambda$等于率失真曲线斜率的绝对值，即多一比特能减少的失真；码率越低，曲线越陡，所以粗量化应配较大的$\lambda$。

DVC把上述手工模块逐个换成网络。光流网络是\ref{sec:vision-video-warping}节的SPyNet，以预训练权重初始化，在率失真目标下一起微调。运动编码器借用学习式图像压缩的自编码器结构：四层$3\times3$、步长2、128通道的卷积，前三层后各接广义除法归一化（定义见\ref{sec:vision-restoration-quantization}节），把$H\times W\times2$的光流变成$H/16\times W/16\times128$的$\bm{m}_t$；运动解码器以三层带逆归一化的转置卷积和一层输出2通道的转置卷积对称地恢复光流。运动补偿网络把取样后的帧、参考帧和解码光流拼成8通道输入，先经$3\times3$、64通道的卷积，再经两次池化和两次上采样组成的编码—解码结构，其中共有6个预激活残差块，下采样一侧的特征经跳跃连接送到上采样一侧，最后由两层$3\times3$卷积输出3通道的预测帧。残差编码器、解码器和估计码率用的概率模型沿用尺度超先验方法（\ref{sec:vision-restoration-hyperprior}节）。表\ref{tab:vision-video-dvc-layers}汇总这些结构；原文报告整个模型约有1100万个参数。

\begin{table*}[tbp]
  \centering
  \small
  \begin{tabularx}{\linewidth}{@{}>{\raggedright\arraybackslash}p{22mm}>{\raggedright\arraybackslash}X>{\raggedright\arraybackslash}X>{\raggedright\arraybackslash}p{17mm}>{\raggedright\arraybackslash}X>{\raggedright\arraybackslash}X@{}}
    \toprule
    阶段或层组 & 层类型 & 卷积核与步长 & 输出通道 & 重复次数 & 输出尺寸\\
    \midrule
    输入 & 当前帧与上一重建帧 & — & 3 & 2 & $H\times W$\\
    光流估计（SPyNet） & 金字塔逐级残差卷积 & $7\times7$，步长1 & 32, 64, 32, 16, 2 & 每级5层，共5级 & $H\times W\times2$\\
    运动编码器 & 卷积，前三层后接广义除法归一化 & $3\times3$，步长2 & 128 & 4 & $H/16\times W/16\times128$\\
    量化 & 取整；训练时加均匀噪声 & — & 128 & 1 & $H/16\times W/16\times128$\\
    运动解码器 & 转置卷积，前三层后接逆归一化 & $3\times3$，步长2 & 128，末层2 & 4 & $H\times W\times2$\\
    运动补偿：入口 & 拼接与卷积 & $3\times3$，步长1 & 64 & 1 & $H\times W\times64$\\
    运动补偿：主体 & 预激活残差块；池化与上采样 & 论文未给出块内配置 & 64 & 残差块6个，池化与上采样各2次 & 回到$H\times W$\\
    运动补偿：出口 & 卷积 & $3\times3$，步长1 & 64，3 & 2 & $H\times W\times3$\\
    残差编码器与解码器 & 尺度超先验的非线性变换 & 论文未给出 & 论文未给出 & 论文未给出 & 与原图同尺寸的残差\\
    码率估计 & 尺度超先验的概率模型 & 论文未给出 & — & — & 每个符号的概率\\
    \bottomrule
  \end{tabularx}
  \caption{DVC的结构，按原论文正文、运动编码器图和补充材料中的运动补偿网络图整理。光流网络的层配置取自SPyNet原论文，DVC只说明以其预训练权重初始化。残差编码器和概率模型沿用尺度超先验方法，DVC未给出所用的通道数；池化的类型与跳跃连接的合并方式，图中也没有标明。原文报告全模型约1100万个参数。}
  \label{tab:vision-video-dvc-layers}
\end{table*}

量化与熵编码沿用\ref{sec:vision-restoration-quantization}节和\ref{sec:vision-restoration-entropy-coding}节的做法：训练时以加均匀噪声代替取整以便求梯度，使用时真正取整为整数符号，再按共享的概率模型熵编码。DVC用学习的概率模型估计两类符号的负对数概率，作为码率$R_{\mathrm{motion},t}$与$R_{\mathrm{residual},t}$；残差的概率模型带有尺度超先验，其辅助信息的比特也计入残差码率，计算方式见式\eqref{eq:vision-hyperprior-rate}。接收端先把比特还原成相同符号，再运行运动和残差解码网络。潜在张量的元素数不等于码流大小，实际长度还受概率分布和编码格式影响。

与混合编码逐块用式\eqref{eq:vision-video-lagrangian}选择模式不同，DVC把率失真权衡放进训练目标。对一帧使用
\begin{equation}
  L_t=\lambda D(\bm{Y}_t,\widehat{\bm{Y}}_t)
       +R_{\mathrm{motion},t}+R_{\mathrm{residual},t}\text{，}
  \label{eq:vision-video-rate-distortion}
\end{equation}
其中$D$在原实现中取均方误差，$\lambda>0$控制质量的权重。把式\eqref{eq:vision-video-lagrangian}两边同除以$\lambda$可见，两者是同一个拉格朗日代价，只是乘子写在失真一侧，所以这里$\lambda$越大，失真受到的惩罚越重，网络倾向于多花比特。原文分别取$\lambda=256,512,1024,2048$训练四个模型，得到四个码率工作点；编码时不再逐块搜索，权衡已经固化在网络参数中。若$D$按像素平均、$R$按每帧比特，$\lambda$的尺度与改用比特/像素时不同，比较时应固定归一化。光流网络也受这个目标约束，因而它学到的运动要兼顾可压缩性与预测用途，不应仅以光流端点误差评价它。

以教学数值说明联合目标的必要性。假设同一帧的方案A用20比特描述运动、80比特描述残差，方案B用10比特描述运动、105比特描述残差，总数分别为100与115。运动分支省了10比特，整个系统却多花15比特，因为粗运动把更多误差留给了残差。这个对照不假设两方案失真相同；若失真不同，还要一起报告。

原文的消融实验给出了联合训练的实际效果。若在整个训练中固定预训练的光流网络，编码光流需要0.044比特/像素，取样所得预测帧的PSNR为27.33 dB；联合训练后只需0.029比特/像素，PSNR反而升到28.17 dB，即学到的运动既更省比特，又更利于预测。去掉运动自编码器、直接量化原始光流，光流要花0.20比特/像素，预测帧PSNR降到约24.3 dB，整体码率—质量曲线在同等码率下下降超过2 dB；去掉运动补偿网络、只用取样结果作预测，同等码率下PSNR约降1.0 dB。在其报告的$\lambda=512,1024,2048$三个工作点上，运动信息约占总码率的36\%、35\%和27\%，质量要求越高，占比越低\scite{vision-dvc2019}。

训练还要处理参考依赖：第$t$帧的输入依赖第$t-1$帧的重建，而后者又依赖更早的帧。DVC采用在线更新的缓存：把每个训练片段上一步的重建帧存入缓存，下一步直接读取作参考，避免每次从序列首帧回算整条链；原文报告，用原始前一帧代替缓存中的重建帧训练，同等码率下约损失0.2 dB。训练数据是Vimeo-90K的短片段，$256\times256$裁剪，批量4，Adam优化，学习率从$10^{-4}$开始在损失停滞时降为十分之一，在两块Titan X上约训练7天。缓存只是训练组织方式；真正编码时仍要顺着当前参数形成的重建链运行，独立解码器不能依赖训练缓存或原视频。

两条路线在比特分配上的差别可以归结为三点。第一，运动的表示：混合编码用块和运动矢量，划分本身要花比特，但描述紧凑；DVC用逐像素光流再压缩成潜在表示，没有块边界，比特由网络自行分配。第二，权衡发生的时刻：混合编码在编码时对每块搜索模式，可以随内容和码率目标调整；DVC在训练时固化权衡，一个$\lambda$对应一个模型，编码时只做一次前向计算。第三，变换的适应性：混合编码的整数变换固定，DVC的非线性变换从数据学到。DVC因此可以联合优化各模块，代价是缺少逐块自适应的模式选择，码率点也只有训练过的几个。

完整视频应统计首帧或关键帧、各帧运动、残差以及实际传输的尺寸、模式等辅助信息。估计码率用于训练，真实码流字节数用于核验传输成本。模型若预先部署在两端，需说明模型版本；若传输中更新模型或概率参数，那些字节也属于必要成本。

\subsection{视频帧的解码重建}
\label{sec:vision-video-decoding}

第三步在接收端按依赖顺序还原每一帧。输入只有码流和事先共享的模型，输出是重建序列。难点在于任何一帧的解码都依赖它的参考，参考出错或缺失会影响后续所有帧，而用户可能想从中间开始观看。

混合编码的解码器按码流顺序解析每块的模式、运动矢量和系数，反量化、反变换后加回预测，再做环路滤波。块状处理在块边界上留下的不连续是最显眼的压缩伪影。H.264规定了自适应的环路\term{去块滤波}（deblocking filter）\scite{vision-wiegand2003h264}：在块边界两侧比较样本差异，若边界处差异明显而又小于随量化参数变化的阈值，就认为它是量化造成的块效应并加以平滑；若差异大到量化的粗细解释不了，就认为是图像中的真实边缘而不处理。因为滤波发生在预测环路之内，滤波后的帧也是后续帧的参考，改善的质量可以传给后续预测；概述报告，在相同客观质量下它通常节省5\%到10\%的码率。HEVC在去块之后又加入\term{样点自适应偏移}（sample adaptive offset, SAO），用少量编码器统计出的参数，以查表方式修正样本幅值\scite{vision-sullivan2012hevc}。

DVC的独立解码器先恢复图像组的首帧，再按参考关系读取后续帧的运动与残差符号，用解码网络得到$\widehat{\bm{v}}_t$和$\widehat{\bm{E}}_t$，执行运动补偿与式\eqref{eq:vision-video-dvc-reconstruction}的相加，然后把当前结果写回参考缓存。原文认为逐像素的运动补偿不产生块效应，因而不需要手工设计的环路滤波和样点自适应偏移；按原文的设计，取样结果中的伪影由运动补偿网络去除，它承担了类似的修正职责。所有运算的尺寸、裁剪、量化和数值约定都应与编码器的本地解码一致。这里出现了混合编码已经解决、学习式编码重新面对的问题：神经网络在不同硬件上的浮点运算结果可能有末位差异，两端的参考于是不完全相同，正是H.264用精确整数反变换避免的漂移，需要在部署时检查。只有从保存的码流独立重建成功，才能说明输出是一套可用的编码结果。

单参考链不能从任意中间帧直接解码，随机访问要从所在图像组的首帧开始恢复依赖。图像组越短，独立帧越多，码率越高，但随机访问等待和错误传播范围越小。DVC的原实验在不同测试集上分别使用长度12和10的图像组。传输丢包还可能破坏后续参考，不能只检查每帧单独的解码函数。多参考或双向预测会改变依赖、缓冲、解码顺序和延迟条件，应作为不同编码配置比较。

DVC原文报告，在其测试设置下，按PSNR和MS-SSIM计算的码率—质量曲线在多数数据集上优于其所配置的H.264编码器，按MS-SSIM与H.265相近或更好\scite{vision-dvc2019}。这一比较对应原文的具体配置：H.264和H.265由FFmpeg的very fast预设生成，图像组长度如上；补充材料按序列给出了相对H.264的BD码率，例如在BasketballDrive序列上，按PSNR计算H.265为$-44.37\%$，DVC为$-23.17\%$。编码标准名不能替代具体编码器、预设、图像组长度和速率控制配置，更不能据此给出当前持续有效的排名。速度也取决于实现：在$352\times288$的视频上，原文测得DVC编码约24.5帧每秒、解码约41帧每秒（单块1080Ti），官方参考软件JM和HM的编码速度分别为2.4和0.35帧每秒，而高度优化的x264和x265分别达到250和42帧每秒。

DVC与BasicVSR都依赖光流和反向取样，比较二者能看清编码问题的特殊之处。在BasicVSR中，光流在使用端由输入帧随时计算，不需付出比特，可以尽量精确；在DVC中，光流必须编码传输，精度要用比特换取，运动网络因此学习可压缩的运动。BasicVSR的参考是低质量输入，误差不会沿参考链放大；DVC的参考是自己的重建，每帧的编码误差都会成为下一帧的起点。前者以质量为唯一目标，后者以率失真为目标，评价时也必须同时报告码率。

\section{处理效果的核验}
\label{sec:vision-video-restoration-evaluation}

一个结果可能逐帧更清楚，却在播放时闪烁；另一个可能节省比特，却要等待整段视频或改变字符。本节要解决的问题是：在什么条件下，两份视频处理结果的比较能说明方法本身的差别，以及“处理成功”需要哪些证据。\ref{sec:vision-video-restoration-challenges}节已经给出了指标的定义，这里关心的是怎样使用它们。核验有三个维度，分别针对不同的失败来源：比较协议排除输入和信息条件不同造成的假差异；内容与时间连续性检查单帧指标看不到的时序错误；质量、码率与延迟则把收益与成本放在一起，按用途选择。

\subsection{比较协议的核验}
\label{sec:vision-video-restoration-protocol}

比较前要确认双方面对的是同一个问题。超分需要固定倍率、下采样与模糊方式、裁边和色彩空间，BasicVSR的两种退化就要分开报告；插帧需要固定两端间隔、目标时刻以及是否以先前生成帧作输入；压缩需要固定输入分辨率、帧率、图像组长度、参考结构和码流统计方式。可见帧范围也属于协议：一个方法看整段，另一个只看过去，单帧质量差异就同时包含信息条件差异。训练与测试划分要按源视频进行，相邻帧或近重复片段落入不同划分会泄漏信息。

指标的汇总方式同样属于协议。PSNR是逐帧计算再平均，还是先平均均方误差再换算，按\ref{sec:vision-video-restoration-challenges}节的算例可以相差数个分贝；只在亮度通道上计算还是在RGB上计算，也会改变数值，BasicVSR的原文表格就对不同数据集采用了不同通道。码率按式\eqref{eq:vision-video-bitrate}由完整码流字节数计算，不能用潜在张量的元素数或训练时的估计码率代替。

计算延迟、等待后帧的延迟和显示缓冲要按式\eqref{eq:vision-video-latency}分别测量。吞吐量描述每秒处理多少帧，延迟描述某帧从拍摄到输出等待多久；报告时还需给出硬件、数值精度、批量、输入长度和峰值内存，避免只比较网络某个局部模块的耗时。DVC与x264的速度比较就说明了这一点：同一种编码思路，参考软件与高度优化的商业实现可以相差两个数量级，速度数字只对具体实现有效。

\subsection{内容与时间连续性的核验}
\label{sec:vision-video-temporal-quality}

协议统一后，要检查输出内容是否正确、播放是否连续。单帧指标PSNR、SSIM和LPIPS仍然适用，但它们逐帧计算再平均，看不到闪烁、重影、拖尾和字符随时间变化。直接惩罚相邻输出之差$\widehat{\bm{X}}_t-\widehat{\bm{X}}_{t-1}$也不行，因为正常运动时相邻帧本来就不同，这样会鼓励画面停滞；时间比较应补偿真实运动，并排除遮挡、新显露和越界区域，这正是式\eqref{eq:vision-video-temporal-error}的扭曲误差的设计。

使用扭曲误差时要说明三件事：对应和掩码来自哪里，掩码怎样得到（例如由式\eqref{eq:vision-video-fb-consistency}的前后向一致性），以及按什么范数和归一化计算。若对应来自被评价方法自己的光流，误差会被方法本身掩盖；光流不准、曝光变化或掩码漏判也会影响它。扭曲误差和tOF回答的问题不同：前者问相邻输出在运动补偿后是否一致，后者问输出的运动是否与真值的运动一致，需要配对真值。两者应与逐帧指标一起报告，因为任何一个单独都可以被投机：模糊输出的扭曲误差低；逐帧补出彼此不一致细节的输出，逐帧感知距离可能较低，播放时却闪烁。

时间差小也不保证内容正确。每帧都输出同一张模糊灰图，可以很稳定却丢失全部细节；每帧都把车牌同一字符改错，也能保持时间一致。因此要同时检查与参考内容的关系及实际用途：逐帧读取车牌，回查原观测，比较识别置信度和字符串随时间的变化，能识别“单帧好看、播放时字符反复变换”的失败。有配对清晰视频时，可比较输出与真值的运动补偿差异；没有真值时，主观连续性、下游任务和原始证据检查更重要，但不能把这些替代检查说成已证明像素恢复正确。

发现失败后，应追溯到具体环节：取样位置错误（对齐）、错误内容被赋予高权重（聚合）、错误沿状态扩散（传播）、两侧都未见的区域由先验补出（插帧合成），或参考链失配（编码）。本章各方法在这些环节上的设计不同，同一种表面现象可能来自不同原因：同样是残影，在运动补偿时域滤波中多半来自过大的时间权重，在RVRT中可能来自错误特征被逐片段传播，在DVC中则可能来自两端参考不一致造成的漂移。

\subsection{质量、码率与延迟的核验}
\label{sec:vision-video-quality-rate-latency}

最后要在收益与成本之间作出选择。单个工作点无法比较两种方法：一个码率更低、质量也更低的结果，可能只是选了不同的$\lambda$或量化参数。正确做法是固定源序列，用多个配置分别形成质量—码率曲线和质量—等待曲线，在相同码率下比较质量，或在相同质量下比较码率；式\eqref{eq:vision-bd-rate}的BD码率把这种比较概括为一个数，但它只在两条曲线重叠的质量区间内有效，曲线交叉时一个平均数会掩盖“低码率时更好、高码率时更差”的情况，应同时给出曲线。若一个配置更省码率却必须等待整段视频，另一个略差但能即时输出，选择取决于用途：观看容忍适量细节损失，车牌读取更在意字符，测量还要求原始观测和时间坐标可追溯。

表\ref{tab:vision-video-tradeoff}汇总不同处理职责最容易遗漏的条件，说明同一“清晰视频”目标为何需要多项核验。比较质量、码率与延迟时，应同时保留对应、可见性、真实细节和解码依赖的检查。

\begin{table}[htbp]
  \centering
  \begin{tabularx}{\linewidth}{@{}lXXX@{}}
    \toprule
    处理职责 & 主要新增信息 & 必须同时检查 & 主要成本\\
    \midrule
    多帧复原 & 邻帧互补内容 & 对应、可见性与错误传播 & 等待范围与缓存\\
    空间超分 & 位移形成的采样差异 & 真实细节与先验生成细节 & 计算量与等待\\
    时间插帧 & 两端的时间约束 & 后帧等待、遮挡与运动歧义 & 后帧等待\\
    帧间压缩 & 已解码参考的可预测内容 & 码流、参考一致与随机访问 & 码率、重排等待与解码依赖\\
    \bottomrule
  \end{tabularx}
  \caption{视频处理的不同收益依赖不同信息条件。比较应固定输入与用途，把右侧两列与单帧质量一起检查，不能据某一项改善推断整个系统已满足要求。}
  \label{tab:vision-video-tradeoff}
\end{table}
\FloatBarrier

\section{本章小结}
\label{sec:vision-video-restoration-summary}

开篇的问题是怎样利用其他时刻的路牌观测得到可读、可放大、可补帧、可传输的视频。本章先把这个问题拆成五项挑战：对齐误差带来伪影，遮挡使邻帧缺少所需内容，时间一致性与逐帧质量相互牵制，读取范围受等待约束，编码还要在运动与残差之间分配比特。回答都围绕帧间对应展开。邻帧只有在对应正确、内容可见时才有用：反向取样决定到哪里读取，前后向一致性筛掉不可见位置，权重决定读来的信息怎样参与，状态传播决定哪些历史仍能影响后续输出。经典的运动补偿时域滤波用按局部差异调整的时间权重做递归平均，其方差公式把降噪强度与残影持续时间联系在一起；VBM4D沿轨迹把相似图块分组协同收缩；RVRT以片段内联合处理和光流引导的多位置注意力，在逐帧循环与整段并行之间取得折中，并把手工的相似度和阈值换成学到的注意力。

输出网格决定了结果中有多少是观测、多少是估计。原网格恢复每个位置都有退化观测；BasicVSR用双向传播取得整段的亚像素观测，再由像素重排扩展空间网格，新增像素依赖邻帧的间接观测；FILM从未观测的中间时刻出发向两侧取样，输出完全是估计。编码则要求共享信息能在解码端重现：H.264/HEVC的混合编码以块运动补偿、整数变换、量化参数和上下文自适应熵编码组成，编码器用拉格朗日代价逐块选择模式，编码器内部复制解码环路，精确整数反变换消除了两端漂移；DVC保留这一数据流，把运动估计、运动压缩、运动补偿和残差变换换成网络，用率失真损失联合训练，使学到的运动更便于压缩，代价是权衡固化在模型中，也重新面对浮点运算的两端一致问题。

本章没有解决的部分同样来自这些依赖。双向方法需要等待整段视频，改成在线要重新设计和验证；大运动、遮挡、切镜和真实时变退化仍会使对应失效，错误还可能沿状态和参考链扩散；插帧和超分中由先验补出的细节无法从输出本身区分出来；极低码率下，运动—残差框架能保留的信息有限，学习式编码的模型部署与跨硬件一致也还需逐项检查。因此，画面更锐利、更平滑或码流更小，都不足以单独证明任务完成；改进需要在相同输入条件、可见范围和用途下，同时核验逐帧质量、时间一致性、码率与延迟。

\paragraph{扩展阅读}
跨帧对齐与传播方向，可以对照本章以光流取样为主的路线阅读几项工作。TDAN用可变形卷积在特征层面直接学习采样偏移，不先显式估计光流；EDVR在滑动窗口内用金字塔级联的可变形对齐和时空注意力融合处理邻帧，是循环方法普及前窗口路线的代表；BasicVSR++在BasicVSR的双向传播上加入二阶网格传播和光流引导的可变形对齐，也就是本章与RVRT对比时提到的后续版本\scite{vision-tdan2020,vision-edvr2019,vision-basicvsrpp2022}VRT是RVRT作者此前提出的整段并行模型，用时间互注意力在帧间交换信息，可用来具体检查并行路线的融合收益与长序列内存代价；FastDVDnet则用两级级联的U形网络直接读取五帧，不做显式运动估计，与依赖对应的经典滤波和VBM4D形成对照\scite{vision-vrt2024,vision-fastdvdnet2020}

插帧方向，Super SloMo把两帧间的双向光流线性组合成中间时刻的近似光流，并预测可见性图处理遮挡，可以一次插出任意多个中间帧；DAIN在光流投影中加入深度，让较近的物体在遮挡处优先；RIFE用一个网络直接估计中间时刻的光流，不再从两帧间光流缩放，换取实时速度\scite{vision-superslomo2018,vision-dain2019,vision-rife2022}三者都可与FILM比较对大运动和去遮挡区域的处理方式。编码方向，H.266/VVC的综述介绍了HEVC之后的混合编码标准怎样加入更灵活的块划分和更多预测工具，并扩展到屏幕内容、全景视频等应用；学习式编码中，FVC把运动估计、压缩与补偿从像素空间移到特征空间，DCVC则不再从预测帧中减出残差，而把时间上下文特征作为编码、解码与熵模型的条件，二者分别从运动表示和残差形式上改变了DVC的数据流\scite{vision-vvc2021,vision-fvc2021,vision-dcvc2021}
